\chapter{相互作用传播与相关函数}\label{chap:interacting-field-propagation}

自由场已经给出多粒子空间、正频谱与局域对易关系。现在让这些自由度发生相互作用，就需要回答局域激发怎样传播，以及不同场插入之间怎样形成相关。运动算子的逆还不足以回答这个问题，必须同时指定真空、时间排序和边界条件。

二点传播子建立后，时间有序演化将局域相互作用送入同一个多粒子空间；Wick 配对把自由场乘积化成可计算的二点函数。对具有稳定外部单粒子极点的过程，外腿约化再从完整相关函数中提取散射振幅。振幅怎样变成探测计数、截面或衰变率，还要结合\cref{chap:relativistic-scattering}的相空间与通量。

\section{传播子与相互作用展开}

Fock 空间给出了自由粒子的态，但一次局域激发怎样从 $x$ 传播到 $y$ 仍由运动方程的逆核决定。逆核还必须带有边界条件，否则同一微分算符有不止一种 Green 函数。选定自由传播以后，相互作用只需在其上作 Dyson 展开；自由场的高阶时序乘积再由 Wick 配对化成二点函数。实验散射态位于遥远的入射区和出射区，因而最后还要从这些相关函数中剥离外腿的自由极点，留下真正的散射振幅。

\subsection{传播子的边界条件}

自由场的二点时序关联函数与运动算符的 Feynman 逆核是同一对象的两种刻画。时序定义先固定真空关联的物理边界条件，随后作用运动算符才显示其 Green 函数性质。先记
\(D^+(x-y)\equiv\langle0|\varphi(x)\varphi(y)|0\rangle\)，它是只含正频质量壳分量的真空二点函数，通常称为 Wightman 函数。对自由实标量场，时间排序只是按两个时刻的先后次序选取 \(D^+(x-y)\) 或 \(D^+(y-x)\)：
\begin{equation}
{
\Delta_F(x-y)
\equiv\langle0|\Tord\{\varphi(x)\varphi(y)\}|0\rangle
=\Theta(x^0-y^0)D^+(x-y)+\Theta(y^0-x^0)D^+(y-x).}
\label{eq:feynman-from-wightman-r3}
\end{equation}
两支 Wightman 函数分别承载正、负频谱分量，阶跃函数按照时间先后选择相应分量。

同一个自由标量二次核可以有不同的 Green 逆算符，因为极点处必须补充边界条件。推迟逆算符在 $x^0<y^0$ 时为零，适合经典初值响应；超前逆算符选择相反因果方向；Feynman 逆算符则编码真空时间排序。对有源方程
\begin{equation*}
\left(\Box+\frac{m^2c^2}{\hbar^2}\right)\varphi(x)=J(x),
\end{equation*}
推迟解为
\begin{equation*}
\varphi(x)=\varphi_{\rm hom}(x)+\int\dd^4 y\,G_R(x-y)J(y),
\end{equation*}
因此 $G_R$ 与 $\Delta_F$ 的物理用途不同，即使它们反演的是同一个微分算符。

Feynman 边界条件在动量空间中由 $\ii0^+$ 规定指定：正能极点位于实轴下方，负能极点位于实轴上方。沿用在 SI 场归一化下，最终得到
\begin{equation}
{
\Delta_F(x-y)
=\int\frac{\dd^4 p}{(2\pi\hbar)^4}
\frac{\ii\hbar^3c}{p^2-m^2c^2+\ii0^+}
 \ee^{-ip\cdot(x-y)/\hbar}.}
\label{eq:scalar-Fourier-r2}
\end{equation}
其极点选择恰好重现\eqnrefs{eq:feynman-from-wightman-r3}：$x^0>y^0$ 时围道只包围正能极点，$x^0<y^0$ 时只包围负能极点。

传播子之所以也是 Green 函数，可以直接由正则对易子看出。作用自由标量逆算符后，非等时区域的 Wightman 函数都满足齐次方程，唯一剩下的项来自时间排序阶跃函数在 $x^0=y^0$ 处的导数。因此


\begin{proposition}[自由标量 Green 核的算符方程]
设质量与 Fourier 约定固定，采用 Feynman 边界值处方；等式在缓增分布意义成立。则有
\begin{equation}
{
\left(\Box_x+\frac{m^2c^2}{\hbar^2}\right)\Delta_F(x-y)
=-\ii\hbar c\,\delta^{(4)}(x-y).}
\label{eq:scalar-green-exact-r2}
\end{equation}
\end{proposition}
\begin{proof}
记 $\tau=x^0-y^0$、$\Omega_{\vb{p}}=E_{\vb{p}}/c$。\eqnrefs{eq:scalar-Fourier-r2} 中的 $p^0$ 积分为

\begin{equation*}
I(\tau)=\int\frac{\dd p^0}{2\pi\hbar}
\frac{\ii\hbar^3c\,\ee^{-ip^0\tau/\hbar}}
{(p^0-\Omega_{\vb{p}}+\ii0^+)(p^0+\Omega_{\vb{p}}-\ii0^+)}.
\end{equation*}
当 $\tau>0$ 时指数要求在下半平面闭合，顺时针积分轮廓给
\begin{equation*}
I(\tau>0)=\frac{\hbar^2c^2}{2E_{\vb{p}}}
 \ee^{-iE_{\vb{p}}\tau/(\hbar c)}.
\end{equation*}
当 $\tau<0$ 时在上半平面闭合，得到
\begin{equation*}
I(\tau<0)=\frac{\hbar^2c^2}{2E_{\vb{p}}}
 \ee^{+iE_{\vb{p}}\tau/(\hbar c)}.
\end{equation*}
恢复空间 Fourier 相位，并在第二支作 $\vb{p}\to-\vb{p}$，就得到\eqnrefs{eq:feynman-from-wightman-r3}。

再从时间序形式直接作用微分算符。第一次对 $x^0$ 求导时会出现
\begin{equation*}
\delta(\tau)[D^+(x-y)-D^+(y-x)],
\end{equation*}
但等时场对易子为零，所以该项消失。第二次求导产生等时 $[\dot\varphi(\vb{x},t),\varphi(\vb{y},t)]$。由正则动量 $\pi=\dot\varphi/c^2$ 与
\begin{equation*}
[\varphi(\vb{x},t),\pi(\vb{y},t)]=\ii\hbar\delta^{(3)}(\vb{x}-\vb{y})
\end{equation*}
得到接触项 $-\ii\hbar c\delta(x^0-y^0)\delta^{(3)}(\vb{x}-\vb{y})$；而 $x\ne y$ 时两支 Wightman 函数的齐次自由标量方程项彼此为零。于是得到\eqnrefs{eq:scalar-green-exact-r2}。
\end{proof}

定义 $G_F=\ii\Delta_F/(\hbar c)$ 后便有 $(\Box+m^2c^2/\hbar^2)G_F=\delta^{(4)}$。因此真空时序关联与 Feynman Green 逆核由同一等时正则代数和同一极点边界条件固定。

\begin{figure}[htbp]
\centering
\begin{tikzfig}[>=Latex]
\draw[lecturearrow] (-4,0)--(4,0) node[right] {$\operatorname{Re}p^0$};
\draw[lecturearrow] (0,-2)--(0,2) node[above] {$\operatorname{Im}p^0$};
\fill (2.2,-0.45) circle (2pt) node[below right] {$+E_{\vb{p}}/c-\ii0^+$};
\fill (-2.2,0.45) circle (2pt) node[above left] {$-E_{\vb{p}}/c+\ii0^+$};
\draw[dashed] (2.2,0)--(2.2,-0.45);
\draw[dashed] (-2.2,0)--(-2.2,0.45);
\end{tikzfig}

\caption{Feynman $\ii0^+$ 规定：正能极点位于下半平面，负能极点位于上半平面；围道的闭合方向由 $x^0-y^0$ 的符号决定。}
\label{fig:feynman-poles-r2}
\end{figure}

\subsection{对量子历史求和与相关函数的源方法}

算符量子力学还有一种完全等价的演化语言：固定初态与末态后，把所有可能的中间历史都作为积分变量，每条历史按 $\ee^{\ii S/\hbar}$ 的相位权重贡献振幅。这种“对历史求和”的表示称为\emph{路径积分}。演化算符的乘法性质、完备关系和受控的时间切片极限直接构造出这一“对历史求和”的表示。场论中的函数积分同样先用格点、模截断或其它调节定义为有限自由度积分，再取连续极限；连续 Green 核由调节后的二次算符逆在该极限中定义。

正则量子化使用 Hilbert 空间和算符代数描述动力学；路径积分则从同一个演化算符出发，把有限时间传播分解成许多短时间传播。取一维原型位置坐标 $\xi$，哈密顿量为
\begin{equation*}
 \hat H=\frac{\hat p^2}{2m_{\rm qm}}+V(\hat\xi),
\end{equation*}
其中 $m_{\rm qm}$ 是原型粒子的质量，SI 单位为 $\mathrm{kg}$；$\hat\xi$ 是位置算符，SI 单位为 $\mathrm m$；与它正则共轭的 $\hat p$ 是一维物理动量算符，满足 $[\hat\xi,\hat p]=\ii\hbar$。把总时间 $T=t_f-t_i$ 分成 $N$ 个小区间
\begin{equation*}
 \Delta t\equiv\frac{T}{N},
 \qquad [\Delta t]=\mathrm s,
\end{equation*}
并在每一小步之间插入位置完备关系
\begin{equation*}
 \id()=\int\dd \xi_n\,|\xi_n\rangle\langle \xi_n|.
\end{equation*}
连续插入 $N-1$ 次后，有限时间核严格分解为
\begin{align}
 \langle \xi_f|\ee^{-\ii \hat H T/\hbar}|\xi_i\rangle
 &=\int\prod_{n=1}^{N-1}\dd \xi_n
 \prod_{n=0}^{N-1}
 \langle \xi_{n+1}|\ee^{-\ii \hat H\Delta t/\hbar}|\xi_n\rangle.
 \label{eq:path-slicing-r15}
\end{align}
这一分解由演化算符的群性质与位置本征态的完备关系给出；单个短时核的构造还需要显式处理动能算符与势能算符的非对易性。

为了计算单个短时核，记
\[
\hat T=\frac{\hat p^2}{2m_{\rm qm}},\qquad \hat V=V(\hat\xi).
\]
在有限格点或有限能量截断 $\Pi_\Lambda$ 下，Lie--Trotter 分裂给
\begin{equation}
 \ee^{-\ii\Delta t(\hat T+\hat V)/\hbar}
 =\ee^{-\ii\Delta t\hat T/\hbar}
  \ee^{-\ii\Delta t\hat V/\hbar}
 +\order((\Delta t)^2),
 \label{eq:trotter-short-r15}
\end{equation}
而单片首个误差算符正比于 $[\hat T,\hat V]$。固定总时间 $T=N\Delta t$ 时，一个有意义的受控量例如
\begin{equation}
 \epsilon_{\rm split}(\Lambda)
 \sim\frac{T\Delta t}{\hbar^2}
 \bigl\|\Pi_\Lambda[\hat T,\hat V]\Pi_\Lambda\bigr\|\ll1.
 \label{eq:trotter-control-regulated-dp73}
\end{equation}
因此“$\Delta t$ 很小”必须相对于实际的不对易能标理解。对未截断的无界算符，一般不能用全空间算符范数写同一估计；若 $\hat T$ 与 $\hat V$ 满足相应的自伴与半有界条件，应使用 Trotter 乘积公式的强极限
\[
 \ee^{-\ii T(\hat T+\hat V)/\hbar}
 =\operatorname*{s-lim}_{N\to\infty}
 \left(\ee^{-\ii T\hat T/(N\hbar)}\ee^{-\ii T\hat V/(N\hbar)}\right)^N.
\]

在动能与势能因子之间插入动量完备关系，短时核为
\begin{align}
 K_{\Delta t}(\xi_{n+1},\xi_n)
 &=\int\frac{\dd p_n}{2\pi\hbar}
 \exp\left\{\frac{\ii}{\hbar}
 \left[p_n(\xi_{n+1}-\xi_n)
 -\Delta t\frac{p_n^2}{2m_{\rm qm}}
 -\Delta t\,V(\xi_n)\right]\right\}
 +\order((\Delta t)^2),
 \label{eq:short-kernel-momentum-r15}\\
 &=\sqrt{\frac{m_{\rm qm}}{2\pi\ii\hbar\Delta t}}
 \exp\left\{\frac{\ii\Delta t}{\hbar}
 \left[
 \frac{m_{\rm qm}}{2}
 \left(\frac{\xi_{n+1}-\xi_n}{\Delta t}\right)^2
 -V(\xi_n)
 \right]\right\}
 +\order((\Delta t)^2).
 \label{eq:short-kernel-config-r15}
\end{align}
第一行显示短时相位是 $p\Delta\xi-H\Delta t$；第二行只是对标准 Fresnel 二次积分的计算结果。

把 $N$ 个短时核相乘，指数中的黎曼和满足
\begin{align}
 &\sum_{n=0}^{N-1}\Delta t
 \left[
 \frac{m_{\rm qm}}{2}
 \left(\frac{\xi_{n+1}-\xi_n}{\Delta t}\right)^2
 -V(\xi_n)
 \right]\\
 &\qquad\xrightarrow[N\to\infty]{}
 \int_{t_i}^{t_f}\dd t\,
 \left[\frac{m_{\rm qm}}{2}\dot\xi^2-V(\xi)\right]
 \equiv S[\xi].
 \label{eq:lattice-action-limit-r15}
\end{align}
同时所有 Fresnel 前因子与中间位置积分共同定义泛函测度
\begin{equation*}
 \mathcal D\xi\equiv
 \lim_{N\to\infty}
 \left(\frac{m_{\rm qm}}{2\pi\ii\hbar\Delta t}\right)^{N/2}
 \prod_{n=1}^{N-1}\dd \xi_n.
\end{equation*}
因此
\begin{equation}
 {
 \langle \xi_f|\ee^{-\ii\hat H T/\hbar}|\xi_i\rangle
 =\int_{\xi(t_i)=\xi_i}^{\xi(t_f)=\xi_f}\mathcal D\xi\,
 \exp\left[\frac{\ii}{\hbar}S[\xi]\right].}
 \label{eq:qm-path-integral-r15}
\end{equation}
\eqnrefs{eq:qm-path-integral-r15} 因而是有限维时间演化在连续插入完备关系并取 Trotter 极限后的连续表示。泛函测度由离散时间切片上的普通多重积分定义；场论中同样先以有限格点或其它调节化离散化定义测度，再取连续极限。

\begin{exbox}[自由粒子传播核]
\eqnrefs{eq:qm-path-integral-r15} 给出任意势能下的路径积分定义，但“对所有路径求和”仍然比较抽象。取 $V(\xi)=0$ 和 $T=t_f-t_i>0$，求传播核并检验短时极限与两段传播的拼接关系。
\textbf{解。} 自由粒子的传播核为
\begin{equation}
 {
 K_0(\xi_f,t_f;\xi_i,t_i)
 =\sqrt{\frac{m_{\rm qm}}{2\pi\ii\hbar T}}
 \exp\!\left[
 \frac{\ii m_{\rm qm}(\xi_f-\xi_i)^2}{2\hbar T}
 \right],
 \qquad T\equiv t_f-t_i>0.}
 \label{eq:free-particle-kernel-dp61}
\end{equation}
这个结果把路径积分里的三个对象同时具体化。指数恰好是连接两个端点的经典直线路径作用量
\begin{equation*}
 S_{\rm cl}
 =\frac{m_{\rm qm}(\xi_f-\xi_i)^2}{2T},
\end{equation*}
前因子则不是可随意省略的“归一化常数”：它负责在 $T\to0^+$ 时恢复 $\delta(\xi_f-\xi_i)$，并保证连续两段传播满足
\begin{equation}
 {
 \int_{-\infty}^{\infty}\dd\xi\,
 K_0(\xi_f,t_f;\xi,t)
 K_0(\xi,t;\xi_i,t_i)
 =K_0(\xi_f,t_f;\xi_i,t_i).}
 \label{eq:free-kernel-composition-dp61}
\end{equation}
因此“路径积分测度”“经典作用量”和“量子演化的群性质”不是三个互不相干的约定；自由粒子已经把它们锁定在同一个可直接检验的闭式中。
\end{exbox}

% -----------------------------------------------------------------------------
% -----------------------------------------------------------------------------

\begin{exbox}[谐振子传播核]
自由粒子核只含动能；真正检验“二次作用量可以精确高斯积分”的下一步，是一维谐振子。它的拉格朗日量为
\begin{equation}
 {L=\frac{m_{\rm qm}}2\dot\xi^2
 -\frac{m_{\rm qm}\omega_0^2}{2}\xi^2,}
 \label{eq:harmonic-lagrangian-dp68}
\end{equation}
其中 $m_{\rm qm}>0$ 是粒子质量，$\omega_0>0$ 是谐振角频率。固定两个端点和传播时间，求经典轨道、经典作用量与量子传播核，并确定前因子和焦散相位。
\textbf{解。} 对 $T=t_f-t_i>0$，固定端点 $\xi(0)=\xi_i$、$\xi(T)=\xi_f$ 的经典轨道为
\begin{equation}
 {
 \xi_{\rm cl}(t)
 =\frac{\xi_i\sin[\omega_0(T-t)]
 +\xi_f\sin(\omega_0t)}{\sin(\omega_0T)}.}
 \label{eq:harmonic-classical-path-dp64}
\end{equation}
先取 $0<T<\pi/\omega_0$ 避开第一焦散点。沿这条经典轨道的作用量是
\begin{equation}
 {
 S_{\rm cl}
 =\frac{m_{\rm qm}\omega_0}{2\sin(\omega_0T)}
 \left[
 (\xi_f^2+\xi_i^2)\cos(\omega_0T)-2\xi_f\xi_i
 \right].}
 \label{eq:harmonic-classical-action-dp64}
\end{equation}
由于作用量对路径是二次泛函，任意路径都可写成经典轨道加零端点涨落，而且两部分在作用量中精确解耦。最终传播核为
\begin{equation}
 {
 K_{\rm ho}(\xi_f,T;\xi_i,0)
 =\sqrt{\frac{m_{\rm qm}\omega_0}
 {2\pi\ii\hbar\sin(\omega_0T)}}
 \exp\!\left[\frac{\ii}{\hbar}S_{\rm cl}(\xi_f,\xi_i;T)\right].}
 \label{eq:harmonic-kernel-dp64}
\end{equation}
这个结果比自由粒子多出的不是某个经验相位，而是两个可以分别追溯的结构：指数来自满足端点条件的\emph{经典}谐振轨道作用量，前因子来自端点固定的量子涨落高斯积分。

先检查 $\omega_0T\ll1$。展开
\begin{equation*}
 \sin(\omega_0T)=\omega_0T+O((\omega_0T)^3),
 \qquad
 \cos(\omega_0T)=1-\frac{\omega_0^2T^2}{2}+\cdots,
\end{equation*}
得到
\begin{equation*}
 S_{\rm cl}
 =\frac{m_{\rm qm}(\xi_f-\xi_i)^2}{2T}
 +O(T),
\end{equation*}
而前因子趋于
\begin{equation*}
 \sqrt{\frac{m_{\rm qm}}{2\pi\ii\hbar T}}.
\end{equation*}
所以\eqnrefs{eq:harmonic-kernel-dp64} 连续回到自由粒子核~\eqnrefs{eq:free-particle-kernel-dp61}。这说明短时间内势能尚未来得及改变主导的局域扩散尺度。

另一个有用的特殊时刻是四分之一周期。取四分之一周期
\begin{equation*}
 T=\frac{\pi}{2\omega_0},
\end{equation*}
则 $\sin(\omega_0T)=1$、$\cos(\omega_0T)=0$，于是
\begin{equation*}
 K_{\rm ho}
 =\sqrt{\frac{m_{\rm qm}\omega_0}{2\pi\ii\hbar}}
 \exp\!\left(-\frac{\ii m_{\rm qm}\omega_0}{\hbar}\xi_f\xi_i\right).
\end{equation*}
这个核除了尺度因子外正是 Fourier 核：谐振子演化四分之一周期把初始坐标分布旋转到共轭动量分布。用相空间语言说，谐振子把 $(\xi,p)$ 作角度 $\omega_0T$ 的刚性旋转；时间透镜、相空间剪切与电子波包聚焦都具有同一种“传播即相空间线性变换”的结构。

再看 $T\to\pi/\omega_0$。此时 $\sin(\omega_0T)\to0$，坐标表示的前因子发散；这不是演化算符失去幺正性，而是所有经典轨道在半周期后把 $\xi_i$ 聚焦到 $-\xi_i$，坐标核趋向带相位的 $\delta(\xi_f+\xi_i)$。跨过这些焦散点以后，平方根分支会积累 Maslov 相位。因此闭\eqnrefs{eq:harmonic-kernel-dp64} 的局部分支由 $T\to0^+$ 的自由核归一化固定，较长时间则通过核的组合律连续延拓。


\medskip

\eqnrefs{eq:harmonic-lagrangian-dp68} 给出二次势中的经典动力学。这个推导只解决一个问题：在给定初末位置后，哪一条经典轨道满足 Euler--Lagrange 方程。端点条件决定两个积分常数，因此不需要任何量子假设。

Euler--Lagrange 方程为
\begin{equation*}
 \ddot\xi_{\rm cl}+\omega_0^2\xi_{\rm cl}=0.
\end{equation*}
一般解写成
\begin{equation*}
 \xi_{\rm cl}(t)=A\cos(\omega_0t)+B\sin(\omega_0t).
\end{equation*}
由 $\xi_{\rm cl}(0)=\xi_i$ 得 $A=\xi_i$；再用 $\xi_{\rm cl}(T)=\xi_f$，
\begin{equation*}
 B=\frac{\xi_f-\xi_i\cos(\omega_0T)}{\sin(\omega_0T)}.
\end{equation*}
代回一般解并利用
\begin{equation*}
 \sin[\omega_0(T-t)]
 =\sin(\omega_0T)\cos(\omega_0t)
 -\cos(\omega_0T)\sin(\omega_0t),
\end{equation*}
即可整理成\eqnrefs{eq:harmonic-classical-path-dp64}。


\medskip

\eqnrefs{eq:harmonic-classical-path-dp64} 已经固定了满足两端边界条件的经典轨道。沿满足运动方程的轨道，谐振子拉格朗日量化成全导数，因此经典作用量由两个端点的 $\xi_{\rm cl}\dot\xi_{\rm cl}$ 直接确定。

利用运动方程，
\begin{align*}
 \frac{\dd}{\dd t}(\xi_{\rm cl}\dot\xi_{\rm cl})
 &=\dot\xi_{\rm cl}^2+\xi_{\rm cl}\ddot\xi_{\rm cl}\\
 &=\dot\xi_{\rm cl}^2-\omega_0^2\xi_{\rm cl}^2.
\end{align*}
所以
\begin{align*}
 S_{\rm cl}
 &=\frac{m_{\rm qm}}2\int_0^T\dd t
 (\dot\xi_{\rm cl}^2-\omega_0^2\xi_{\rm cl}^2)\\
 &=\frac{m_{\rm qm}}2[\xi_{\rm cl}\dot\xi_{\rm cl}]_0^T.
\end{align*}
由\eqnrefs{eq:harmonic-classical-path-dp64}，
\begin{align*}
 \dot\xi_{\rm cl}(0)
 &=\frac{\omega_0}{\sin(\omega_0T)}
 [\xi_f-\xi_i\cos(\omega_0T)],\\
 \dot\xi_{\rm cl}(T)
 &=\frac{\omega_0}{\sin(\omega_0T)}
 [\xi_f\cos(\omega_0T)-\xi_i].
\end{align*}
将 $\xi_{\rm cl}(0)=\xi_i$、$\xi_{\rm cl}(T)=\xi_f$ 与这两个端点速度代入上式，合并同类项后得到\eqnrefs{eq:harmonic-classical-action-dp64}。


\medskip

经典作用量只决定传播核的端点相位；还需要确定所有零端点量子涨落给出的归一化前因子。因为\eqnrefs{eq:harmonic-lagrangian-dp68} 对路径是二次的，将路径写成经典轨道与零端点涨落之和后，交叉项由经典运动方程精确消失，剩余涨落积分只依赖传播时间而不依赖两个端点。

写
\begin{equation*}
 \xi(t)=\xi_{\rm cl}(t)+\delta\xi(t),
 \qquad \delta\xi(0)=\delta\xi(T)=0.
\end{equation*}
作用量展开为
\begin{align*}
 S[\xi]
 =&\,S_{\rm cl}
 +m_{\rm qm}\int_0^T\dd t
 [\dot\xi_{\rm cl}\,\delta\dot\xi-\omega_0^2\xi_{\rm cl}\,\delta\xi]\\
 &+\frac{m_{\rm qm}}2\int_0^T\dd t
 ((\delta\dot\xi)^2-\omega_0^2(\delta\xi)^2).
\end{align*}
中间的线性项分部积分后为
\begin{equation*}
 -m_{\rm qm}\int_0^T\dd t
 (\ddot\xi_{\rm cl}+\omega_0^2\xi_{\rm cl})\,\delta\xi=0,
\end{equation*}
故
\begin{equation*}
 K_{\rm ho}=N(T)\exp\!\left(\frac{\ii}{\hbar}S_{\rm cl}\right).
\end{equation*}
把这一形式代入传播核满足的 Schr\"odinger 方程，端点相关的 $O(\hbar^0)$ 项自动满足 Hamilton--Jacobi 方程，剩余项给
\begin{equation*}
 \frac{\dot N}{N}
 =-\frac{1}{2m_{\rm qm}}
 \frac{\partial^2S_{\rm cl}}{\partial\xi_f^2}
 =-\frac{\omega_0}{2}\cot(\omega_0T).
\end{equation*}
积分得到
\begin{equation*}
 N(T)=\frac{C}{\sqrt{\sin(\omega_0T)}}.
\end{equation*}
短时间极限必须恢复自由粒子核，而 $\sin(\omega_0T)\sim\omega_0T$，因此
\begin{equation*}
 C=\sqrt{\frac{m_{\rm qm}\omega_0}{2\pi\ii\hbar}}.
\end{equation*}
代回 $K_{\rm ho}=N\exp(\ii S_{\rm cl}/\hbar)$ 就得到\eqnrefs{eq:harmonic-kernel-dp64}。
\end{exbox}








场论先把场 $\varphi(x)$ 在格点或有限模基上调节成有限组积分变量，并以调节极限定义连续函数积分。为了用同一个对象一次生成任意阶时间有序相关函数，在作用量中加入一个任意的普通数值函数 $\mathcal J(x)$，让它只以线性形式与场相乘。对这个外加函数求一次泛函导数会插入一个 $\varphi(x)$，求 $n$ 次就会插入 $n$ 个场，因此把“每个给定 $\mathcal J$ 对应的真空函数积分”整体看成 $\mathcal J$ 的泛函最方便。这个用于生成相关函数的对象称为\emph{生成泛函}，定义为
\begin{equation}
 Z[\mathcal J]
 =\int\mathcal D\varphi\,
 \exp\left\{\frac{\ii}{\hbar}
 \left[S[\varphi]+\frac1c\int\dd^4 x\,\mathcal J(x)\varphi(x)\right]\right\}.
 \label{eq:scalar-ZJ-r3}
\end{equation}
源 $\mathcal J(x)$ 的作用只是给每个时空点的场插入一个线性探针。对 $\mathcal J$ 求泛函导数会把一个场算符拉下来：
\begin{equation*}
 \frac{\hbar c}{\ii}\frac{\delta Z[\mathcal J]}{\delta \mathcal J(x)}
 =\int\mathcal D\varphi\,\varphi(x)
 \ee^{\frac{\ii}{\hbar}(S+S_\mathcal J)}.
\end{equation*}
因此归一化后
\begin{equation}
 \left.\left(\frac{\hbar c}{\ii}\right)^n
 \frac{1}{Z[0]}
 \frac{\delta^n Z[\mathcal J]}{\delta \mathcal J(x_1)\cdots\delta \mathcal J(x_n)}
 \right|_{\mathcal J=0}
 =\langle0|\Tord\{\hat\varphi(x_1)\cdots\hat\varphi(x_n)\}|0\rangle.
 \label{eq:ZJ-correlators-r3}
\end{equation}
因此 $Z[\mathcal J]$ 的各阶源泛函导数同时生成全部时序 $n$ 点 Green 函数。

自由场的作用量是二次型，因此生成泛函由高斯积分完全确定。令有限维积分变量为 $\vb{z}\in\mathbb R^n$，二次核为可逆矩阵 $\mathcal K$，源为 $\mathcal J$。在 Feynman $+\ii0^+$ 规定固定平方根支以后，定义
\begin{equation*}
 \mathcal N_{\mathcal K}
 \equiv
 \int\dd^nz\,
 \exp\!\left(\frac{\ii}{2\hbar}\vb{z}^{\mathsf T}\mathcal K\vb{z}\right),
\end{equation*}
则配方给出严格高斯恒等式
\begin{equation}
 \int\dd^nz\,
 \exp\left[\frac{\ii}{2\hbar}\vb{z}^{\mathsf T}\mathcal K\vb{z}+\frac{\ii}{\hbar}\mathcal J^{\mathsf T}\vb{z}\right]
 =
 \mathcal N_{\mathcal K}
 \exp\left[-\frac{\ii}{2\hbar}\mathcal J^{\mathsf T}\mathcal K^{-1}\mathcal J\right],
 \qquad
 \mathcal N_{\mathcal K}=(2\pi\ii\hbar)^{n/2}(\det\mathcal K)^{-1/2}.
 \label{eq:gaussian-finite-r2}
\end{equation}
对自由场，调节后的二次核可直接求逆；连续极限中的逆算符就是带所选边界条件的 Green 核。由于取 $x^0=ct$，$\dd^4 x=c\,\dd t\,\dd^3 x$。实自由标量场作分部积分后写成
\begin{equation*}
 S_0[\varphi]
 =-\frac1{2c}\int\dd^4 x\,\dd^4 y\,
 \varphi(x)\mathcal K(x,y)\varphi(y),
\end{equation*}
其中
\begin{equation*}
 \mathcal K(x,y)
 =\left(\Box_x+\frac{m^2c^2}{\hbar^2}\right)
 \delta^{(4)}(x-y).
\end{equation*}
Feynman 真空边界条件选择 $\mathcal K_F^{-1}$。若源项写成
$S_\mathcal J=(1/c)\int\dd^4 x\,\mathcal J\varphi$，完成平方后所有 $\mathcal J$-依赖都由 $\mathcal K_F^{-1}$ 控制。与 \eqnrefs{eq:scalar-green-exact-r2} 对比，二者的精确关系为
\begin{equation}
 {
 \Delta_F(x,y)=-\ii\hbar c\,\mathcal K_F^{-1}(x,y).}
 \label{eq:scalar-kernel-inverse-r2}
\end{equation}
完成平方可把自由生成泛函写成
\begin{equation}
 {
 \frac{Z_0[\mathcal J]}{Z_0[0]}
 =\exp\left[
 -\frac{1}{2\hbar^2c^2}
 \int\dd^4 x\,\dd^4 y\,
 \mathcal J(x)\Delta_F(x-y)\mathcal J(y)
 \right],}
 \label{eq:free-generating-functional-r3}
\end{equation}
指数中的二次型由自由 Feynman 传播子控制。两次源导数给出一个二点函数，四次源导数给出三种两两配对；Wick 定理的组合结构由此已经包含在高斯生成泛函中。

若作用量写成 $S=S_0+S_{\rm int}$，则
\begin{align}
 Z[\mathcal J]
 &=\int\mathcal D\varphi\,
 \ee^{\frac{\ii}{\hbar}S_0}
 \ee^{\frac{\ii}{\hbar}S_{\rm int}[\varphi]}
 \ee^{\frac{\ii}{\hbar c}\int \mathcal J\varphi},\\
 &=\exp\left[
 \frac{\ii}{\hbar}S_{\rm int}\!\left[
 \frac{\hbar c}{\ii}\frac{\delta}{\delta \mathcal J}
 \right]\right] Z_0[\mathcal J].
 \label{eq:interacting-Z-from-free-r3}
\end{align}
第二行用源泛函导数代替相互作用作用量中的场。展开相互作用指数产生耦合常数的微扰级数，而自由生成泛函的每一对源导数产生一个 $\Delta_F$。因此路径积分中的内部线与算符方法中的 Wick 收缩都来自自由二次核的逆；相互作用顶角则来自 $S_{\rm int}$ 对场的多项式结构。

% -----------------------------------------------------------------------------
% -----------------------------------------------------------------------------

\begin{align}
 \frac{\delta^2\mathcal Z_0}
 {\delta\mathcal J(x_1)\delta\mathcal J(x_2)}
 =&\;\mathcal Z_0
 \frac{\delta\mathscr Q}{\delta\mathcal J(x_1)}
 \frac{\delta\mathscr Q}{\delta\mathcal J(x_2)}
 +\mathcal Z_0
 \frac{\delta^2\mathscr Q}
 {\delta\mathcal J(x_1)\delta\mathcal J(x_2)},\\
 \frac{\delta^2\mathscr Q}
 {\delta\mathcal J(x_1)\delta\mathcal J(x_2)}
 =&-\frac{1}{\hbar^2c^2}\Delta_F(x_1-x_2).
 \label{eq:qft-QJ-second-derivative-dp6}
\end{align}

\begin{proposition}[二次高斯源指数的二阶微分]
设源泛函为正文所定义的对称二次型，核的分布配对有意义；在有限模正则化下取普通导数，再按分布极限理解连续源微分。则有以下结论（原式 \eqnrefs{eq:qft-QJ-second-derivative-dp6}）：
\begin{align*}
 \frac{\delta^2\mathcal Z_0}
 {\delta\mathcal J(x_1)\delta\mathcal J(x_2)}
 =&\;\mathcal Z_0
 \frac{\delta\mathscr Q}{\delta\mathcal J(x_1)}
 \frac{\delta\mathscr Q}{\delta\mathcal J(x_2)}
 +\mathcal Z_0
 \frac{\delta^2\mathscr Q}
 {\delta\mathcal J(x_1)\delta\mathcal J(x_2)},\\
 \frac{\delta^2\mathscr Q}
 {\delta\mathcal J(x_1)\delta\mathcal J(x_2)}
 =&-\frac{1}{\hbar^2c^2}\Delta_F(x_1-x_2).
\end{align*}
\end{proposition}
\begin{proof}
正文\eqnrefs{eq:free-generating-functional-r3} 把自由理论的全部源依赖写成一个二次指数。二次高斯源指数的两次源微分恰好产生一个 Feynman 传播子。这里计算这一二阶导数；四点及更高点则由同一高斯指数的高阶微分产生，并由 Wick 配对统一组织。

记
\begin{equation*}
 \mathscr Q[\mathcal J]
 \equiv
 -\frac{1}{2\hbar^2c^2}
 \int\dd^4 x\,\dd^4 y\,
 \mathcal J(x)\Delta_F(x-y)\mathcal J(y),
 \qquad
 \mathcal Z_0[\mathcal J]=\ee^{\mathscr Q[\mathcal J]}.
\end{equation*}
由于实标量场的 Feynman 核满足
$\Delta_F(x-y)=\Delta_F(y-x)$，第一次泛函导数为
\begin{align*}
 \frac{\delta\mathscr Q}{\delta\mathcal J(x_1)}
 &=-\frac{1}{2\hbar^2c^2}
 \int\dd^4 y\,\Delta_F(x_1-y)\mathcal J(y)
 -\frac{1}{2\hbar^2c^2}
 \int\dd^4 x\,\mathcal J(x)\Delta_F(x-x_1)\\
 &=-\frac{1}{\hbar^2c^2}
 \int\dd^4 y\,\Delta_F(x_1-y)\mathcal J(y).
\end{align*}
于是
\begin{equation*}
 \frac{\delta\mathcal Z_0}{\delta\mathcal J(x_1)}
 =\mathcal Z_0
 \frac{\delta\mathscr Q}{\delta\mathcal J(x_1)}.
\end{equation*}
再对 $\mathcal J(x_2)$ 求导，乘积法则给出
\begin{align*}
 \frac{\delta^2\mathcal Z_0}
 {\delta\mathcal J(x_1)\delta\mathcal J(x_2)}
 =&\;\mathcal Z_0
 \frac{\delta\mathscr Q}{\delta\mathcal J(x_1)}
 \frac{\delta\mathscr Q}{\delta\mathcal J(x_2)}
 +\mathcal Z_0
 \frac{\delta^2\mathscr Q}
 {\delta\mathcal J(x_1)\delta\mathcal J(x_2)},\\
 \frac{\delta^2\mathscr Q}
 {\delta\mathcal J(x_1)\delta\mathcal J(x_2)}
 =&-\frac{1}{\hbar^2c^2}\Delta_F(x_1-x_2),
\end{align*}
这就是正文\eqnrefs{eq:qft-QJ-second-derivative-dp6}。在零源处，两个一阶导数都含显式 $\mathcal J$ 而消失，所以剩下的二阶核直接给出自由二点函数；更高点的配对结构同样来自这一高斯性质，并可整理成 Wick 配对。
\end{proof}




散射与量子修正还需要把相关函数按“能否分成彼此独立的部分”继续分类。先考虑若干外部时空点上的场插入。如果一个贡献能够分解成两个互不相连的因子的乘积，那么其中一组外部点与另一组之间没有任何传播链把它们连接起来；这种贡献称为\emph{非连通}。反之，只要所有外部插入都属于同一个连接网络，就称为\emph{连通}。这个区分由相关函数能否分解为互不连通因子的乘积来定义。$Z[J]$ 同时生成连通和非连通贡献；对真空归一化后的生成泛函取对数，正好把乘积型非连通部分消去，只留下连通累积量。因此写成
\begin{equation}
\mathcal Z[J]\equiv\frac{Z[J]}{Z[0]},
\qquad
W[J]\equiv-\ii\hbar\ln\mathcal Z[J].
\label{eq:connected_generating_functional_definition}
\end{equation}
$W$ 的 SI 量纲为作用量。由于源项采用
$S_J=c^{-1}\int\dd^4 x\,J(x)\varphi(x)$，逐项求导有
\begin{align*}
\frac{\delta W}{\delta J(x)}
&=-\ii\hbar\frac1{\mathcal Z}
\frac{\delta\mathcal Z}{\delta J(x)}
=\frac1c\langle\hat\varphi(x)\rangle_J.
\end{align*}
因此定义含源经典场
\begin{equation}
{
\varphi_c(x)\equiv c\frac{\delta W[J]}{\delta J(x)}.}
\label{eq:qft_classical_field_from_W}
\end{equation}
第二次导数已经显示“取对数”怎样消去非连通乘积。由
\(
W=-\ii\hbar\ln\mathcal Z
\)
先求一次导数，得到
\begin{equation*}
\frac{\delta W}{\delta J_1}
=-\ii\hbar\frac{1}{\mathcal Z}
\frac{\delta\mathcal Z}{\delta J_1}.
\end{equation*}
再求第二次导数：
\begin{align*}
\frac{\delta^2W}{\delta J_1\delta J_2}
=-\ii\hbar\left[
\frac1{\mathcal Z}
\frac{\delta^2\mathcal Z}{\delta J_1\delta J_2}
-\frac1{\mathcal Z^2}
\frac{\delta\mathcal Z}{\delta J_1}
\frac{\delta\mathcal Z}{\delta J_2}
\right].
\end{align*}
方括号的第二项正好减去二点矩中的一阶矩乘积，因此留下二阶连通累积量。第三次求导会进一步产生“完整三点矩减去三个一阶乘二阶乘积，再加两个一阶矩三重乘积”的三阶累积量。继续求导时，每一阶都按乘积法则生成外部点的所有分组方式，而对数产生的交替组合系数恰好逐层减去能够分解成较低阶乘积的部分。因此 $W$ 的任意高阶源导数只保留不能再分解为独立乘积的连通关联函数。

连通相关函数仍然可以包含“由一条内部传播子串联起来的两个较小块”。为了找出不能再由这种单条内部线拼接得到的基本块，再定义\emph{单粒子不可约}（简称 1PI）：一个连通图若切断任意一条内部传播线后仍不会分成两个互不相连的部分，就称为 1PI。一般连通图可以由若干 1PI 顶角块和完整传播子重新拼接，因此 1PI 对象适合作为量子修正后的基本动力学核。

数学上，从“外源 $J$ 是自变量”改到“外源诱导的平均场 $\varphi_c$ 是自变量”需要一次 Legendre 变换。这与经典力学中在不同共轭变量之间变换势函数的做法相同：不是增加新的物理自由度，而是换一种更适合写运动方程的变量。由此定义\emph{1PI 有效作用量}
\begin{equation}
{
\Gamma_{\rm eff}[\varphi_c]
\equiv W[J]
-\frac1c\int\dd^4 x\,J(x)\varphi_c(x),}
\label{eq:onepi_effective_action_definition}
\end{equation}
其中右端的 $J$ 要理解为由\eqnrefs{eq:qft_classical_field_from_W} 反解成 $\varphi_c$ 的泛函。$\Gamma_{\rm eff}$ 的 $n$ 次泛函导数生成 1PI 的 $n$ 点顶角核；特别地，二阶导数是完整二点传播子的逆核。变分时
\begin{align*}
\delta\Gamma_{\rm eff}
&=\frac1c\int\dd^4 x\,\varphi_c(x)\delta J(x)
-\frac1c\int\dd^4 x\,[\delta J(x)\varphi_c(x)+J(x)\delta\varphi_c(x)]\\
&=-\frac1c\int\dd^4 x\,J(x)\delta\varphi_c(x),
\end{align*}
故
\begin{equation}
{
\frac{\delta\Gamma_{\rm eff}}{\delta\varphi_c(x)}=-\frac{J(x)}{c}.}
\label{eq:effective_action_eom}
\end{equation}
物理真空或平均场对应 $J=0$，所以精确量子运动方程是
$\delta\Gamma_{\rm eff}/\delta\varphi_c=0$。这条式子说明“量子修正后的运动方程”不是把若干圈图能量修正手工加回经典方程，而是由同一个 $\Gamma_{\rm eff}$ 的驻值条件统一生成。

二点层级还给出一个关键逆关系。由
$\delta\varphi_c(x)/\delta J(y)=c\,\delta^2W/\delta J(x)\delta J(y)$，以及
$\delta J(y)/\delta\varphi_c(z)=-c\,\delta^2\Gamma_{\rm eff}/\delta\varphi_c(y)\delta\varphi_c(z)$，链式法则
\[
\int\dd^4 y\,
\frac{\delta\varphi_c(x)}{\delta J(y)}
\frac{\delta J(y)}{\delta\varphi_c(z)}
=\delta^{(4)}(x-z)
\]
给
\begin{equation}
\int\dd^4 y\,
W^{(2)}(x,y)\,\Gamma_{\rm eff}^{(2)}(y,z)
=-\frac1{c^2}\delta^{(4)}(x-z).
\label{eq:connected_onepi_inverse_relation}
\end{equation}
所以精确连通传播子与精确 1PI 二点核互为逆算符。把完整逆传播子与自由逆传播子之差组织成一个二点 1PI 核，这个核称为\emph{自能}：它汇总粒子在传播途中发生并最终回到同一粒子通道的全部不可约虚过程。把自由逆传播子与自能重新组合就得到 Dyson 关系 $G^{-1}=G_0^{-1}-\Sigma$；QED 的电子自能和光子真空极化都是这一二点结构的具体实现。

在任意背景 $\varphi_c$ 附近把场分解为
$\varphi=\varphi_c+\delta\varphi$。作用量 Taylor 展开为
\begin{align*}
S[\varphi_c+\delta\varphi]
&=S[\varphi_c]
+\int\dd^4 x\,S^{(1)}(x)\,\delta\varphi(x)\\
&\quad+\frac12\iint\dd^4 x\dd^4 y\,
\delta\varphi(x)S^{(2)}(x,y)\delta\varphi(y)+O((\delta\varphi)^3).
\end{align*}
在背景满足相应源方程后线性项被抵消，高斯涨落积分产生行列式。于是圈图展开的前两阶具有普适结构
\begin{equation}
{
\Gamma_{\rm eff}[\varphi_c]
=S[\varphi_c]
+\frac{\ii\hbar}{2}\Tr\ln S^{(2)}[\varphi_c]
+O(\hbar^2),}
\label{eq:one_loop_effective_action_general}
\end{equation}
其中 $\Tr$ 同时包含时空积分和内部指标。费米型高斯因 Grassmann 行列式位于分子而多一个负号。\eqnrefs{eq:one_loop_effective_action_general} 说明第一阶量子涨落修正完全由经典背景周围二次涨落算符的泛函行列式决定。这个 $O(\hbar)$ 泛函行列式就是图形展开中的一圈贡献：它表示经典背景周围高斯二次涨落对经典作用量的第一阶量子修正。

Grassmann 高斯积分与源平移恒等式给出费米型生成泛函的代数基础。对具体 Dirac 场，有限维矩阵 $\mathsf M$ 被连续时空与自旋指标上的 Dirac 二次核替代，Grassmann 外源记为 $\mathcal J_\psi(x)$ 与 $\overline{\mathcal J}_\psi(x)$；它们是反对易的外部探针，行列式结构和二次源项仍保持同一代数形式。

费米型场的路径积分需要反对易 Grassmann 变量。有限维 Grassmann 高斯恒等式给出其行列式结构，因此其原型不是 $(\det\mathcal K)^{-1/2}$，而是
\begin{equation}
 \int\prod_i\dd\bar\theta_i\,\dd\theta_i\,
 \ee^{-\bar\theta\,\mathsf M\,\theta}
 =\det\mathsf M,
 \label{eq:grassmann-det-r2}
\end{equation}
带 Grassmann 源时完成平方产生 $\mathsf M^{-1}$。因此对易场与反对易场的二次核都通过逆算符控制二点函数；差别在于普通高斯积分给出行列式的负半次幂，而 Grassmann 高斯积分给出行列式的一次幂。在选择具体粒子表示以前，这一有限维恒等式只规定“反对易变量的二次核由逆矩阵产生二点函数”；具体矩阵核必须由该场的时空表示与作用量另行确定。

% Detailed derivation gathered at the end of the owning structure.



\subsection{时间有序相互作用展开}

相互作用理论的困难可以说得很具体：自由哈密顿量 $\hat H_0$ 已经能够精确求解，而完整哈密顿量 $\hat H=\hat H_0+\hat H_{\rm int}$ 一般不能。若完全改用 Schr\"odinger 绘景，所有时间依赖都压在态上，自由演化与相互作用演化混在一起；若完全改用 Heisenberg 绘景，算符又按照完整 $\hat H$ 演化，同样难以求解。相互作用绘景的目的，就是把这两部分有意拆开：算符只承担已知的自由运动，态只承担待展开的相互作用。这样，微扰参数真正只出现在 $\hat H_{\rm int}$ 中。

把哈密顿量分为
\begin{equation*}
 \hat H=\hat H_0+\hat H_{\rm int}.
\end{equation*}
Schr\"odinger 绘景中态满足
\begin{equation*}
 \ii\hbar\partial_t|\Psi(t)\rangle_S
 =\hat H|\Psi(t)\rangle_S.
\end{equation*}
定义相互作用表象
\begin{equation}
 |\Psi(t)\rangle_I
 =\ee^{+\ii\hat H_0t/\hbar}|\Psi(t)\rangle_S,
 \qquad
 \hat O_I(t)=\ee^{+\ii\hat H_0t/\hbar}\hat O_S \ee^{-\ii\hat H_0t/\hbar}.
 \label{eq:interaction-picture-def-r3}
\end{equation}
由这一分工得到相互作用绘景哈密顿量
\begin{equation}
 \hat H_I(t)
 =\ee^{+\ii\hat H_0t/\hbar}
 \hat H_{\rm int}
 \ee^{-\ii\hat H_0t/\hbar},
 \label{eq:HI-interaction-r3}
\end{equation}
而态只满足
\begin{equation}
 {
 \ii\hbar\frac{\dd}{\dd t}|\Psi(t)\rangle_I
 =\hat H_I(t)|\Psi(t)\rangle_I.}
 \label{eq:interaction-sch-r2}
\end{equation}
因此自由场的平面波展开、产生湮灭代数和传播规律都可以继续使用；微扰真正需要处理的只有 $\hat H_I(t)$。

相互作用为什么会改变自由粒子数，可以先在算符层直接看见。对实标量场的四次相互作用，
\begin{equation*}
 \hat\varphi^4\sim(\hat a+\hat a^\dagger)^4,
\end{equation*}
其中同时含有 $a^{\dagger4}$、$a^{\dagger3}a$、$a^{\dagger2}a^2$、$a^\dagger a^3$ 和 $a^4$ 等结构，所以 $\hat H_{\rm int}$ 会连接不同的自由 Fock 扇区。某一项能否成为真实散射或衰变过程，还要再满足四动量守恒和外腿质量壳条件；“算符允许改变粒子数”与“运动学允许该过程发生”是两个不同条件。

定义相互作用绘景演化算符
\begin{equation*}
 |\Psi(t)\rangle_I
 =\hat U_I(t,t_0)|\Psi(t_0)\rangle_I,
\end{equation*}
则
\begin{equation}
 \ii\hbar\partial_t\hat U_I(t,t_0)
 =\hat H_I(t)\hat U_I(t,t_0),
 \qquad
 \hat U_I(t_0,t_0)=\id().
 \label{eq:UI-diff-r3}
\end{equation}
不同时间的 $\hat H_I$ 一般不对易，因此演化算符由微分方程~\eqnrefs{eq:UI-diff-r3} 的 Volterra 积分形式及其逐次迭代定义；第 $n$ 阶系数具有一般形式
\begin{equation}
 \hat U_I^{(n)}(t,t_0)
 =\left(-\frac{\ii}{\hbar}\right)^n
 \int_{t_0}^{t}\dd t_1
 \int_{t_0}^{t_1}\dd t_2\cdots
 \int_{t_0}^{t_{n-1}}\dd t_n\,
 \hat H_I(t_1)\cdots\hat H_I(t_n).
 \label{eq:dyson-general-simplex-dp73}
\end{equation}
积分域已经对任意 $n$ 强制 $t_1>t_2>\cdots>t_n$。把超立方体的 $n!$ 个时间排序区域合并，得到精确的形式解


\begin{proposition}[有序指数的 Volterra 展开]
设相互作用图景中的生成元有界且时间可积，或已在有限维正则化中定义；迭代 Volterra 方程得到收敛 Dyson 级数，未正则化场论按形式级数使用。则有
\begin{equation}
 {
 \hat U_I(t,t_0)
 =\Tord\exp\left[-\frac{\ii}{\hbar}
 \int_{t_0}^{t}\dd t'\,\hat H_I(t')\right].}
 \label{eq:dyson-r2}
\end{equation}
\end{proposition}
\begin{proof}
由\eqnrefs{eq:UI-diff-r3} 积分得到
\begin{equation*}
 \hat U_I(t,t_0)=\id()-\frac{\ii}{\hbar}
 \int_{t_0}^{t}\dd t_1\,\hat H_I(t_1)\hat U_I(t_1,t_0).
\end{equation*}
从 $\hat U_I^{[0]}=\id()$ 开始，把右端的 $\hat U_I$ 反复代回同一个 Volterra 积分。第 $N+1$ 次代回写成
\begin{equation*}
 \hat U_I^{[N+1]}(t,t_0)
 =\id()-\frac{\ii}{\hbar}\int_{t_0}^{t}\dd t_1\,
 \hat H_I(t_1)\hat U_I^{[N]}(t_1,t_0).
\end{equation*}
对 $N$ 作归纳：若第 $n$ 阶齐次项具有正文\eqnrefs{eq:dyson-general-simplex-dp73} 的有序单纯形形式，则再代回一次 Volterra 核只是在最外层增加 $t_0<t_{n+1}<t_n$ 的积分，因此对任意 $n\ge1$
\begin{equation*}
 \hat U_I^{(n)}(t,t_0)
 =\left(-\frac{\ii}{\hbar}\right)^n
 \int_{t_0<t_n<\cdots<t_1<t}\dd t_1\cdots\dd t_n\,
 \hat H_I(t_1)\cdots\hat H_I(t_n).
\end{equation*}
除去等时零测集，超立方体 $[t_0,t]^n$ 被 $n!$ 个时间排列单纯形分割；时间排序算符在每个单纯形上恢复同一有序乘积，所以
\begin{equation*}
 \hat U_I^{(n)}
 =\frac1{n!}\left(-\frac{\ii}{\hbar}\right)^n
 \int_{t_0}^{t}\dd t_1\cdots\dd t_n\,
 \Tord\{\hat H_I(t_1)\cdots\hat H_I(t_n)\}.
\end{equation*}
从 $n=0$ 起求和即得\eqnrefs{eq:dyson-r2}。若任意两时刻的 $H_I$ 彼此对易，时间排序可去掉，结果退化为普通指数；否则普通指数并不满足\eqnrefs{eq:UI-diff-r3}。
\end{proof}

这种由非交换时间演化产生的有序展开通常称为 Dyson 级数。时间排序由\eqnrefs{eq:UI-diff-r3} 的一般解唯一确定，并在图展开中固定各相互作用算符的次序。二阶只是\eqnrefs{eq:dyson-general-simplex-dp73} 的第一个非平凡检验：
\begin{align}
 \hat U_I(t,t_0)
 =&\;\id()-\frac{\ii}{\hbar}\int_{t_0}^{t}\dd t_1\,\hat H_I(t_1)\nonumber\\
 &+\left(-\frac{\ii}{\hbar}\right)^2
 \int_{t_0}^{t}\dd t_1\int_{t_0}^{t_1}\dd t_2\,
 \hat H_I(t_1)\hat H_I(t_2)+\cdots .
 \label{eq:dyson-ordered-domain-r2}
\end{align}

有限 $t_i,t_f$ 的演化由 \eqnrefs{eq:dyson-r2} 定义。散射实验比较的是“相互作用发生前已经彼此远离的自由入射态”和“相互作用结束后再次彼此远离的自由出射态”。把每个自由入射态映射成相应自由出射态的线性算符称为\emph{散射算符}，通常记为 $\hat S$，也称 $S$ 矩阵。要从有限时间演化得到它，必须规定相互作用怎样在遥远过去和未来关闭；先用缓慢开关
\begin{equation*}
 \hat H_I(t)\longrightarrow \ee^{-\epsilon|t|}\hat H_I(t),
 \qquad \epsilon>0,
\end{equation*}
定义
\begin{equation}
 \hat S
 =\lim_{\epsilon\to0^+}
 \Tord\exp\left[-\frac{\ii}{\hbar}
 \int_{-\infty}^{+\infty}\dd t\,
 \ee^{-\epsilon|t|}\hat H_I(t)\right].
 \label{eq:S-adiabatic-r3}
\end{equation}
于是
\begin{equation}
 S_{fi}=\langle f|\hat S|i\rangle,
 \qquad
 \hat S=\id()+\ii\hat T_S.
 \label{eq:S-T-def-r3}
\end{equation}
$\id()$ 对应没有发生散射，$\hat T_S$ 是无量纲的散射跃迁算符。对时间无关散射，还可以抽出能量守恒 $\delta$ 函数，并定义具有能量量纲的定态 $\hat T(E)$。Lippmann--Schwinger 预解式把这一能量分辨的 $\hat T(E)$ 与同一精确散射映射联系起来；Born 级数、Fermi 黄金律和 Feynman 展开只是针对不同条件的下游展开。


% -----------------------------------------------------------------------------
% Exact scattering foundations before perturbative Feynman rules
% -----------------------------------------------------------------------------

有限时间演化算符描述任意给定时间窗内的动力学；散射算符则进一步要求入射波包在遥远过去、出射波包在遥远未来都已经离开相互作用区。这个渐近条件把“器件内发生了什么”压缩成自由通道之间的映射，也使 Born 近似、Dyson 截断和 Feynman 图都能追溯到同一个精确散射对象。

设
\begin{equation}
 \hat H=\hat H_0+\hat V,
 \label{eq:H-scattering-r5}
\end{equation}
其中 $\hat H_0$ 的绝对连续谱描述远离相互作用区的自由通道。M{\o}ller 波算符定义为
\begin{align}
 \hat\Omega_+
 &\equiv
 \operatorname*{s-lim}_{t\to-\infty}
 \ee^{+\ii\hat Ht/\hbar}
 \ee^{-\ii\hat H_0t/\hbar}
 \hat P_{\rm ac}(\hat H_0),\\
 \hat\Omega_-
 &\equiv
 \operatorname*{s-lim}_{t\to+\infty}
 \ee^{+\ii\hat Ht/\hbar}
 \ee^{-\ii\hat H_0t/\hbar}
 \hat P_{\rm ac}(\hat H_0).
 \label{eq:moller-wave-operators-r5}
\end{align}
若这些强极限存在，则
\begin{equation*}
 |\psi^{(+)}\rangle=\hat\Omega_+|\psi_{\rm in}\rangle,
 \qquad
 |\psi^{(-)}\rangle=\hat\Omega_-|\psi_{\rm out}\rangle,
\end{equation*}
而散射算符定义为
\begin{equation}
 \boxed{\hat S=\hat\Omega_-^\dagger\hat\Omega_+.}
 \label{eq:exact-S-waveop-r5}
\end{equation}
它把遥远过去的自由入射通道映射到遥远未来的自由出射通道。Born 级数、Dyson 展开和 Feynman 图都是求这个精确散射映射的具体展开方法。

定态散射把同一问题改写到固定能量 $E$ 上。设 $|\phi_E\rangle$ 是自由 Hamilton 量 $\hat H_0$ 的广义本征态。由于 $E$ 位于连续谱上，$E-\hat H_0$ 不能作普通有界逆；入射边界条件由预解式边界值
\begin{equation}
 \hat G_0^{(+)}(E)
 \equiv\lim_{\epsilon\downarrow0}
 \frac{1}{E-\hat H_0+\ii\epsilon}
 \label{eq:G0-plus-r5}
\end{equation}
选定。完整入射散射态满足
\begin{equation}
 \boxed{
 |\psi_E^{(+)}\rangle
 =
 |\phi_E\rangle
 +\hat G_0^{(+)}(E)\hat V
 |\psi_E^{(+)}\rangle.}
 \label{eq:LS-state-r5}
\end{equation}

\begin{proposition}[预解式构造的散射边界值方程]
设 $H=H_0+V$，相关预解式满足极限吸收原则；广义本征态方程在选定加权检验空间中成立。则有以下结论（原式 \eqnrefs{eq:LS-state-r5}）：
\begin{equation*}
 \boxed{
 |\psi_E^{(+)}\rangle
 =
 |\phi_E\rangle
 +\hat G_0^{(+)}(E)\hat V
 |\psi_E^{(+)}\rangle.}
\end{equation*}
\end{proposition}
\begin{proof}
连续谱能量 $E$ 上不能把 $E-H_0$ 当作普通有界算符直接求逆，因此正文\eqnrefs{eq:G0-plus-r5} 先用 $+\ii0$ 选定出射预解式边界值。把定态 Schrödinger 方程移项后作用这一边界值，并保留自由齐次解，就得到正文\eqnrefs{eq:LS-state-r5}。



令
\begin{equation*}
 (\hat H_0+\hat V)|\psi_E^{(+)}\rangle
 =E|\psi_E^{(+)}\rangle.
\end{equation*}
移项得到
\begin{equation*}
 (E-\hat H_0)|\psi_E^{(+)}\rangle
 =\hat V|\psi_E^{(+)}\rangle.
\end{equation*}
因为 $E$ 位于 $\hat H_0$ 的连续谱上，$E-\hat H_0$ 不能作为普通有界算符直接求逆；入射边界条件由预解式的边界值
\begin{equation*}
 \hat G_0^{(+)}(E)
 \equiv
 \lim_{\epsilon\downarrow0}
 \frac{1}{E-\hat H_0+\ii\epsilon}
 \end{equation*}
选定。于是
\begin{equation*}
 {
 |\psi_E^{(+)}\rangle
 =
 |\phi_E\rangle
 +\hat G_0^{(+)}(E)\hat V
 |\psi_E^{(+)}\rangle,}
 \end{equation*}
其中 $(E-\hat H_0)|\phi_E\rangle=0$。若要求出射边界条件，则改用
\begin{equation*}
 \hat G_0^{(-)}(E)
 =
 \lim_{\epsilon\downarrow0}
 \frac{1}{E-\hat H_0-\ii\epsilon}.
\end{equation*}
所以散射边界条件与 Green 函数在连续谱两侧的解析边界值是同一结构。
\end{proof}

第一项保留指定的自由入射波，第二项是相互作用产生并按出射边界传播的散射波。

把相互作用真正产生的源项定义为
\begin{equation}
 \hat T(E)|\phi_E\rangle
 \equiv
 \hat V|\psi_E^{(+)}\rangle.
 \label{eq:T-def-r5}
\end{equation}
则精确跃迁算符满足
\begin{equation}
 \boxed{
 \hat T(E)
 =
 \hat V+\hat V\hat G_0^{(+)}(E)\hat T(E).}
 \label{eq:T-exact-r5}
\end{equation}
$\hat T(E)$ 因而是把完整散射态压缩到自由通道上的能量分辨算符；Born 级数由这一精确方程对 $\hat V$ 逐次迭代得到。

波算符把自由通道与完整散射态联系起来，同时保持能量生成元的作用：
\begin{equation}
 \boxed{\hat H\hat\Omega_\pm=\hat\Omega_\pm\hat H_0.}
 \label{eq:moller-intertwining-r68}
\end{equation}

\begin{proposition}[波算符的时间演化交织]
设 $H,H_0$ 自伴，正文的强波算符极限存在；生成元恒等式只在自由生成元的定义域上陈述。则有以下结论（原式 \eqnrefs{eq:moller-intertwining-r68}）：
\begin{equation*}
 \boxed{\hat H\hat\Omega_\pm=\hat\Omega_\pm\hat H_0.}
\end{equation*}
\end{proposition}
\begin{proof}
正文\eqnrefs{eq:moller-wave-operators-r5} 用遥远过去和未来的强极限定义 Møller 波算符。把完整时间演化移过这一极限并重命名极限变量，可以证明波算符把 $H_0$ 的时间平移交织成 $H$ 的时间平移；在零时间处求导便得到正文\eqnrefs{eq:moller-intertwining-r68}。



由定义考虑
\begin{align*}
 \ee^{-\ii\hat H\tau/\hbar}\hat\Omega_+
 &=
 \lim_{t\to-\infty}
 \ee^{-\ii\hat H\tau/\hbar}
 \ee^{+\ii\hat Ht/\hbar}
 \ee^{-\ii\hat H_0t/\hbar},\\
 &=
 \lim_{t\to-\infty}
 \ee^{+\ii\hat H(t-\tau)/\hbar}
 \ee^{-\ii\hat H_0(t-\tau)/\hbar}
 \ee^{-\ii\hat H_0\tau/\hbar},\\
 &=\hat\Omega_+\ee^{-\ii\hat H_0\tau/\hbar}.
\end{align*}
在 $\tau=0$ 对两边求导，
\begin{equation*}
 {
 \hat H\hat\Omega_+
 =\hat\Omega_+\hat H_0.}
 \end{equation*}
对 $\hat\Omega_-$ 需要把极限改为 $t\to+\infty$。对任意有限 $\tau$，
\begin{align*}
 \ee^{-\ii\hat H\tau/\hbar}\hat\Omega_-
 &=\lim_{t\to+\infty}
 \ee^{-\ii\hat H\tau/\hbar}
 \ee^{+\ii\hat Ht/\hbar}
 \ee^{-\ii\hat H_0t/\hbar},\\
 &=\lim_{t\to+\infty}
 \ee^{+\ii\hat H(t-\tau)/\hbar}
 \ee^{-\ii\hat H_0(t-\tau)/\hbar}
 \ee^{-\ii\hat H_0\tau/\hbar},\\
 &=\hat\Omega_-\ee^{-\ii\hat H_0\tau/\hbar}.
\end{align*}
在 $\tau=0$ 处分别对两侧求导得到
\[
 -\frac{\ii}{\hbar}\hat H\hat\Omega_-
 =-\frac{\ii}{\hbar}\hat\Omega_-\hat H_0,
\]
因此
\[
 \hat H\hat\Omega_-=\hat\Omega_-\hat H_0.
\]
于是若 $|\phi_E\rangle$ 是 $\hat H_0$ 的广义能量本征态，
\[
 \hat H\,\hat\Omega_\pm|\phi_E\rangle
 =\hat\Omega_\pm\hat H_0|\phi_E\rangle
 =E\,\hat\Omega_\pm|\phi_E\rangle,
\]
所以 $\hat\Omega_\pm|\phi_E\rangle$ 是完整哈密顿量在同一连续能量上的入、出散射态。

上述微分应作用于 $\phi\in\Dom(H_0)$：有界 $\Omega_\pm$ 将自由强差商极限送到完整演化的强差商极限。由 Stone 生成元域判据，$\Omega_\pm\phi\in\Dom(H)$ 且 $H\Omega_\pm\phi=\Omega_\pm H_0\phi$，从而避免对任意 Hilbert 向量形式求无界导数。
\end{proof}

因此 $\hat\Omega_\pm|\phi_E\rangle$ 与 $|\phi_E\rangle$ 具有相同连续谱能量。对精确 $T$ 方程逐次迭代，在余项于所研究子空间趋零时得到 Born 展开
\begin{equation}
 \boxed{
 \hat T(E)=\hat V+\hat V\hat G_0^{(+)}\hat V
 +\hat V\hat G_0^{(+)}\hat V\hat G_0^{(+)}\hat V+\cdots.}
 \label{eq:born-series-r68}
\end{equation}

\begin{proposition}[预解式的 Neumann 级数]
在所论 Banach 空间中，假设 $G_0(z)V$ 有界且 $\|G_0(z)V\|<1$；实能轴上须另具相应极限吸收控制。则有以下结论（原式 \eqnrefs{eq:born-series-r68}）：
\begin{equation*}
 \boxed{
 \hat T(E)=\hat V+\hat V\hat G_0^{(+)}\hat V
 +\hat V\hat G_0^{(+)}\hat V\hat G_0^{(+)}\hat V+\cdots.}
\end{equation*}
\end{proposition}
\begin{proof}
正文\eqnrefs{eq:T-exact-r5} 已经是完整散射的精确算符方程。将右端的 $T$ 连续代回自身，可把一次、两次、三次散射按 $V G_0^{(+)}$ 的次数排列；只有当剩余项在所研究子空间上可忽略时，有限恒等式才变成正文\eqnrefs{eq:born-series-r68} 的 Born 展开。


从精确方程
\begin{equation*}
T=V+VG_0^{(+)}T
\end{equation*}
出发。第一次把右端的 $T$ 再代回：
\begin{align*}
T
&=V+VG_0^{(+)}\big(V+VG_0^{(+)}T\big)\\
&=V+VG_0^{(+)}V+VG_0^{(+)}VG_0^{(+)}T.
\end{align*}
再迭代一次最后的 $T$：
\begin{align*}
T
={}&V+VG_0^{(+)}V+VG_0^{(+)}VG_0^{(+)}V\\
&+VG_0^{(+)}VG_0^{(+)}VG_0^{(+)}T.
\end{align*}
重复这一过程 $N$ 次得到有限阶恒等式
\begin{equation*}
T=\sum_{n=0}^{N}(VG_0^{(+)})^nV+(VG_0^{(+)})^{N+1}T.
\end{equation*}
若余项在所研究初末态或工作子空间上趋零，才可写成 Born 级数
\begin{equation*}
T=V+VG_0^{(+)}V+VG_0^{(+)}VG_0^{(+)}V+\cdots.
\end{equation*}
因此把一阶、二阶 Born 截断分别定义为
\begin{equation*}
T^{[1]}\equiv V,
\qquad
T^{[2]}\equiv V+VG_0^{(+)}V.
\end{equation*}
上面的有限阶恒等式给出相应精确余项
\begin{equation*}
T-T^{[1]}=VG_0^{(+)}T,
\qquad
T-T^{[2]}=(VG_0^{(+)})^2T.
\end{equation*}
只有这些余项在所研究矩阵元或工作子空间范数中足够小时，$T^{[1]}$ 或 $T^{[2]}$ 才是受控近似。无量纲控制对象因此是“自由传播一次以后再散射”的组合，例如投影到工作子空间后的 $\|G_0^{(+)}V\|$ 或相应矩阵元。

在命题的有界算子条件下，令 $B=G_0(z)V$、$r=\|B\|<1$，有 $\|(I-B)^{-1}-\sum_{n=0}^N B^n\|\le r^{N+1}/(1-r)$。因此该展开具有严格收敛与截断余项；边界值处的形式微扰不能自动继承此范数条件。
\end{proof}

$\hat T$ 的定义由精确 Lippmann--Schwinger 方程给出；Born 近似只发生在级数截断时，其控制量是自由预解式与势的组合，而不是孤立的 $\hat V$。

对时间无关散射，时间平移不变性把散射算符限制在同一能量壳上。采用前文固定的连续谱归一化，$S$ 与能量分辨 $T$ 有如下精确关系。
\begin{proposition}[定态散射的能量壳关系]
设 $H=H_0+V$ 自伴且时间无关，入、出波算符在所论散射空间存在，定态出射预解式边界值与入射散射态存在。假设相互作用图景的绝热极限与下述波包矩阵元积分相容；按前文的连续谱态及 $T(E)$ 归一化，以下关系在能量分布意义下成立：
\begin{equation}
 \boxed{
 \langle f|\hat S|i\rangle
 =\langle f|i\rangle
 -2\pi\ii\,\delta(E_f-E_i)\,
 \langle f|\hat T(E_i)|i\rangle.}
 \label{eq:S-T-energy-r68}
\end{equation}
\end{proposition}
\begin{proof}
这里要分清两件事：\(\delta(E_f-E_i)\) 来自无限时间平移对称，而 \(T(E_i)\) 包含的是作用期间所有重复散射。二者可以从精确相互作用绘景积分方程中直接分离。

设
\begin{equation*}
\hat H_0|i\rangle=E_i|i\rangle,
\qquad
\hat H_0|f\rangle=E_f|f\rangle,
\end{equation*}
并令入射定态散射态
\begin{equation*}
|\psi_i^{(+)}\rangle
\equiv\hat\Omega_+|i\rangle.
\end{equation*}
由前面已证明的波算符交织关系
\(\hat H\hat\Omega_+=\hat\Omega_+\hat H_0\)，有
\begin{equation*}
\hat H|\psi_i^{(+)}\rangle
=E_i|\psi_i^{(+)}\rangle.
\end{equation*}
因此该态在 Schr\"odinger 绘景只积累定态相位，而其相互作用绘景表示为
\begin{align*}
|\psi_{i,I}^{(+)}(t)\rangle
&=\ee^{+\ii\hat H_0t/\hbar}
\ee^{-\ii\hat Ht/\hbar}|\psi_i^{(+)}\rangle\\
&=\ee^{+\ii\hat H_0t/\hbar}
\ee^{-\ii E_it/\hbar}|\psi_i^{(+)}\rangle.
\end{align*}

精确相互作用绘景传播子满足 Volterra 方程
\begin{equation*}
\hat U_I(+\infty,-\infty)
=\id()-\frac{\ii}{\hbar}
\int_{-\infty}^{+\infty}\dd t\,
\hat V_I(t)\hat U_I(t,-\infty).
\end{equation*}
对入射态 \(|i\rangle\)，绝热极限下
\(\hat U_I(t,-\infty)|i\rangle=|\psi_{i,I}^{(+)}(t)\rangle\)。因此
\begin{align*}
\langle f|\hat S|i\rangle
={}&\langle f|i\rangle
-\frac{\ii}{\hbar}
\int_{-\infty}^{+\infty}\dd t\,
\langle f|\hat V_I(t)|\psi_{i,I}^{(+)}(t)\rangle.
\end{align*}
把
\(\hat V_I(t)=\ee^{+\ii\hat H_0t/\hbar}\hat V\ee^{-\ii\hat H_0t/\hbar}\)
和上面的定态形式代入，\(\hat H_0\) 在 bra 上给出 \(E_f\)，于是时间依赖完全提出：
\begin{align*}
\langle f|\hat V_I(t)|\psi_{i,I}^{(+)}(t)\rangle
&=\ee^{+\ii(E_f-E_i)t/\hbar}
\langle f|\hat V|\psi_i^{(+)}\rangle.
\end{align*}

再由 Lippmann--Schwinger 方程
\begin{equation*}
|\psi_i^{(+)}\rangle
=|i\rangle+\hat G_0^{(+)}(E_i)\hat V|\psi_i^{(+)}\rangle
\end{equation*}
以及正文对 \(\hat T(E)\) 的定义，可得
\begin{equation*}
\hat V|\psi_i^{(+)}\rangle
=\hat T(E_i)|i\rangle.
\end{equation*}
这一等式把任意次数的
\(V G_0^{(+)}V G_0^{(+)}\cdots V\)
全部收进了 \(T(E_i)\)，剩余时间积分于是给出整体能量守恒：
\begin{equation*}
\int_{-\infty}^{+\infty}\dd t\,
\ee^{+\ii(E_f-E_i)t/\hbar}
=2\pi\hbar\,\delta(E_f-E_i).
\end{equation*}
于是
\begin{equation*}
\langle f|\hat S|i\rangle
=\langle f|i\rangle
-2\pi\ii\,\delta(E_f-E_i)
\langle f|\hat T(E_i)|i\rangle,
\end{equation*}
即正文\eqnrefs{eq:S-T-energy-r68}。

这个推导同时说明了两个常被混淆的对象：\(\delta(E_f-E_i)\) 只反映时间平移不变性；\(T(E_i)\) 才承载相互作用区内全部多次散射动力学。相对论局域场论再把空间平移对称性并入同一逻辑，得到四维动量守恒 \(\delta^{(4)}\)。
\end{proof}
若散射波算符在所讨论子空间渐近完备，则 $\hat S$ 幺正。写 $\hat S=\id()+\ii\hat T_S$ 后可等价改写为
\begin{equation}
 \boxed{
 -\ii(\hat T_S-\hat T_S^\dagger)
 =\hat T_S^\dagger\hat T_S,
 \qquad
 2\Im\langle i|\hat T_S|i\rangle
 =\sum_X|\langle X|\hat T_S|i\rangle|^2.}
 \label{eq:optical-unitarity-r68}
\end{equation}

\begin{proposition}[散射矩阵的光学恒等式]
设入出波算符渐近完备于同一个散射子空间，使 $S$ 幺正，并取正文 $S,T$ 归一化。则有以下结论（原式 \eqnrefs{eq:optical-unitarity-r68}）：
\begin{equation*}
 \boxed{
 -\ii(\hat T_S-\hat T_S^\dagger)
 =\hat T_S^\dagger\hat T_S,
 \qquad
 2\Im\langle i|\hat T_S|i\rangle
 =\sum_X|\langle X|\hat T_S|i\rangle|^2.}
\end{equation*}
\end{proposition}
\begin{proof}
正文\eqnrefs{eq:exact-S-waveop-r5} 在渐近完备的散射子空间上由两个波算符组成，因此保持总概率。把 $S$ 写成单位算符与跃迁部分之和，展开 $S^\dagger S=\id()$，再在前向入射态之间插入完整出射通道，就得到正文\eqnrefs{eq:optical-unitarity-r68}。



若波算符渐近完备，则散射子空间上
\begin{equation*}
 \hat S^\dagger\hat S=\id().
\end{equation*}
写
\begin{equation*}
 \hat S=\id()+\ii\hat T_S,
\end{equation*}
其中 $\hat T_S$ 是无量纲跃迁算符。代入得
\begin{align*}
 (\id()-\ii\hat T_S^\dagger)
 (\id()+\ii\hat T_S)
 &=\id(),\\
 \ii(\hat T_S-\hat T_S^\dagger)
 +\hat T_S^\dagger\hat T_S
 &=0.
\end{align*}
所以
\begin{equation*}
 {
 -\ii(
 \hat T_S-\hat T_S^\dagger)
 =
 \hat T_S^\dagger\hat T_S.}
 \end{equation*}
在两侧夹入同一入射态并插入完整的在壳出射通道：
\begin{equation*}
 2\,\operatorname{Im}
 \langle i|\hat T_S|i\rangle
 =
 \sum_X
 |\langle X|\hat T_S|i\rangle|^2.
 \end{equation*}
右端是所有开放通道的正概率之和。对二体入射，将连续态求和写成 Lorentz 不变相空间后，便得到光学定理：前向振幅的虚部与总截面成正比。
\end{proof}

第二式是在完整在壳出射通道上插入单位算符后的前向矩阵元形式；加入连续相空间测度后即成为通常的光学定理。

外部粒子由 $t\to\pm\infty$ 的渐近条件选成自由哈密顿量的质量壳本征态，因此外部动量满足各自在壳条件。内部传播来自预解式或时间序二点函数，只受各顶角四动量守恒约束，并没有独立的渐近自由边界条件。因此内部四动量一般满足
\begin{equation*}
 q^2\neq m^2c^2.
\end{equation*}
“虚粒子”只是对这类未被外部质量壳条件固定的内部传播变量的简称；它不是一颗在短时间内“违反能量守恒”的真实粒子。

\begin{note}[$S$ 矩阵与有限时间传播]
无限时间散射可用
\begin{equation*}
 \hat V(t)\longrightarrow
 \ee^{-\epsilon|t|}\hat V,
 \qquad
 \epsilon\downarrow0,
\end{equation*}
选择相互作用的真空和边界条件。对有限长度、有限脉冲的自由电子量子光学器件，更基本的对象仍是有限 $t_i,t_f$ 的传播子 $\hat U(t_f,t_i)$；只有当入、出波包在器件外已经形成清楚的自由通道时，再把有限时间传播压缩成 $S$-矩阵才有意义。
\end{note}



\subsection{自由高斯场的配对结构与外腿约化}

Dyson 展开已经把相互作用问题化成自由场的时序乘积。自由理论是高斯的，因此高阶时序乘积不含新的独立统计信息：它们都可由二点函数的配对组合得到。引入一个紧支撑经典源 $\mathcal J(x)$，可把任意场数的配对一次编码进生成恒等式
\begin{equation}
{
 \Tord\exp\!\left[\frac{\ii}{\hbar c}\int\dd^4 x\,\mathcal J(x)\varphi(x)\right]
 =:\!\exp\!\left[\frac{\ii}{\hbar c}\int\dd^4 x\,\mathcal J(x)\varphi(x)\right]\!:
 \exp\!\left[-\frac{1}{2\hbar^2c^2}\int\dd^4 x\,\dd^4 y\,
 \mathcal J(x)\Delta_F(x-y)\mathcal J(y)\right].}
\label{eq:wick-generating-identity-dp74}
\end{equation}
这里第一个指数保留未配对的正规序场，第二个纯 $c$ 数指数一次产生一个二点收缩。对源作任意阶泛函导数后再令 $\mathcal J=0$，便一次性得到所有场数的配对组合。记
\[
\Delta_{rs}\equiv
\langle0|\Tord\{\varphi_r\varphi_s\}|0\rangle,
\]
并令 $\mathfrak M_n$ 为指标集合 $\{1,\ldots,n\}$ 的全部互不重叠部分配对。若 $V(P)$ 是配对 $P$ 中已经出现的指标集合，则一般恒等式为


\begin{proposition}[自由高斯真空的 Wick 配对]
设态为自由场的零均值高斯真空，$:\ :$ 为相对于该真空的正规序，两点收缩为对应的真空时序核，场乘积按同一时序处方及正则化定义。下式成立；热态或压缩高斯态须改用相应状态的 Wick 正规序与收缩：
\begin{equation}
{
\Tord\!\left\{\prod_{r=1}^{n}\varphi_r\right\}
=\sum_{P\in\mathfrak M_n}
\left[\prod_{(r,s)\in P}\Delta_{rs}\right]
:\!\prod_{k\notin V(P)}\varphi_k\!: .}
\label{eq:wick-general-r3}
\end{equation}
\end{proposition}
\begin{proof}
把生成恒等式写成“正规序线性源”与“收缩二次源”的乘积：
\begin{equation*}
 \mathcal T\exp\!\left[
 \frac{\ii}{\hbar c}\int\dd^4x\,\mathcal J(x)\varphi(x)
 \right]
 =:\exp\!\left[
 \frac{\ii}{\hbar c}\int\dd^4x\,\mathcal J(x)\varphi(x)
 \right]:
 \exp\!\left[-\frac{1}{2\hbar^2c^2}
 \int\dd^4x\dd^4y\,\mathcal J(x)\Delta_F(x-y)\mathcal J(y)
 \right].
\end{equation*}
要得到 $n$ 个指定场 $\varphi_r\equiv\varphi(x_r)$ 的时间序乘积，对两边作 $n$ 次源泛函微分后令 $\mathcal J=0$。非零项中，每一个微分只有两种去处：命中正规序指数的一次源，留下一个未收缩场；或者与另一个微分共同命中二次源指数中的同一个 $\mathcal J(x)\mathcal J(y)$，产生一条二点收缩 $\Delta_F(x_r-x_s)$。

固定恰有 $k$ 条收缩。需要从 $n$ 个标签中选择 $2k$ 个参与收缩，再把它们配成 $k$ 对。带标签配对数为
\begin{equation*}
 \binom{n}{2k}\frac{(2k)!}{2^k k!}
 =\frac{n!}{(n-2k)!2^k k!}.
\end{equation*}
这个组合因子可以直接从二次指数的展开系数看出。其第 $k$ 阶是
\begin{equation*}
 \frac{1}{k!}
 \left[-\frac{1}{2\hbar^2c^2}
 \int\dd^4x\dd^4y\,
 \mathcal J(x)\Delta_F(x-y)\mathcal J(y)
 \right]^k.
\end{equation*}
$n$ 次微分中命中这 $2k$ 个源的排列产生 $(2k)!$；每条二次因子内部交换两个端点产生 $2^k$ 重数，而 $k$ 条二次因子之间交换又产生 $k!$ 重数。原系数中的 $1/(2^kk!)$ 恰好把这些“同一无序配对的重复计数”除掉，所以每一个具体部分配对只留下系数 1。

因此若 $P$ 表示标签集合 $\{1,\ldots,n\}$ 的任意部分配对，$V(P)$ 表示所有已参与配对的标签，则
\begin{equation*}
 \mathcal T\!\left\{\prod_{r=1}^n\varphi_r\right\}
 =\sum_{P}
 \left[\prod_{(r,s)\in P}\Delta_F(x_r-x_s)\right]
 :\!\prod_{j\notin V(P)}\varphi_j\!:,
\end{equation*}
即正文\eqnrefs{eq:wick-general-r3}。

例如 $n=4$ 时，$k=2$ 的完全收缩数为
$4!/(2^2 2!)=3$，正是
$\Delta_{12}\Delta_{34}$、$\Delta_{13}\Delta_{24}$、$\Delta_{14}\Delta_{23}$ 三项。这个低阶检查同时说明 Wick 定理的组合系数来自源导数，而不是费曼图规则额外规定的权重。
\end{proof}

这里空配对给出完整正规序项，完全配对给出纯传播子乘积，其余项包含尚未收缩的正规序场。这个“任意自由高斯时序乘积等于所有互不重叠二点配对之和”的结构通常称为玻色 Wick 定理。Feynman 图中的内部线只是二点缩并的图形化记账，不是经典粒子轨迹。

两场恒等式只是 $n=2$ 的特例：
\begin{equation}
 {
 \Tord\{\varphi(x)\varphi(y)\}
 =:\!\varphi(x)\varphi(y)\!:+\Delta_F(x-y).}
 \label{eq:wick-two-field-r68}
\end{equation}

对四点真空关联，只有完全配对在真空期望值中存活。\eqnrefs{eq:wick-general-r3} 的三个完全配对直接给出
\begin{equation}
 {
 \langle0|\Tord\{\varphi_1\varphi_2\varphi_3\varphi_4\}|0\rangle
 =\Delta_{12}\Delta_{34}
 +\Delta_{13}\Delta_{24}
 +\Delta_{14}\Delta_{23}.}
 \label{eq:wick-four-scalar-example-dp48}
\end{equation}
这个四点公式是一般配对定理的最低非平凡实例；更高点函数由同一配对规则递归生成。

费米场也按所有允许的配对展开，但交换两个 Grassmann 奇算符会产生负号。定义一般自由收缩
\begin{equation*}
 S_{ij}\equiv\langle0|\Tord\{\psi_i\bar\psi_j\}|0\rangle,
\end{equation*}
则完全收缩是带置换奇偶性的配对和；闭合费米子闭合圈因而额外带一个 $-1$。这一负号来自反对易代数本身，而不是 QED 顶角的特殊规定。
% Generic fermionic Gaussian pairing combinatorics before QED specialization.
自由高斯理论的高阶关联函数由二点函数决定。玻色场交换两个算符不会改变符号；费米场每交换两个 Grassmann 奇算符都会产生一个负号。因此一般高阶费米关联函数的每一种完全配对都必须带上把原有算符顺序重排到该配对顺序所需的置换宇称。后面的全同费米子交换号与闭合费米子圈负号都来自这一同一代数结构。


设按原始时间序列排列的 $2n$ 个自由费米型场算符统一记为
\[
\Psi_1,\Psi_2,\ldots,\Psi_{2n},
\]
其中每个 $\Psi_r$ 可以是 $\psi$ 或 $\bar\psi$，并把二点缩并记为
\begin{equation}
 C_{rs}\equiv
 \langle0|\Tord\{\Psi_r\Psi_s\}|0\rangle_{\rm contraction}.
 \label{eq:generic-fermion-contraction-dp8}
\end{equation}
只有场类型与量子数允许的二点缩并才非零。令 $\mathfrak P_{2n}$ 表示集合 $\{1,\ldots,2n\}$ 的全部完全配对；对每个配对 $P$，$(-1)^{\pi(P)}$ 定义为把原始算符序列重排成各配对相邻时所产生的费米置换符号。自由高斯费米场的一般 $2n$ 点函数于是为
\begin{equation}
 \boxed{
 \langle0|\Tord\{\Psi_1\Psi_2\cdots\Psi_{2n}\}|0\rangle
 =\sum_{P\in\mathfrak P_{2n}}
 (-1)^{\pi(P)}
 \prod_{(r,s)\in P} C_{rs}.}
 \label{eq:fermion-wick-general-dp73}
\end{equation}
这一“任意高阶高斯关联函数等于所有二点配对之和，并由费米置换宇称定号”的结果通常称为费米 Wick 定理。其母结构是\eqnrefs{eq:fermion-wick-general-dp73}，低阶公式只是代入后的检验。

例如取 $n=2$ 且算符次序为 $\psi_1\bar\psi_2\psi_3\bar\psi_4$，两个允许的完全收缩具有相反宇称，因此
\begin{equation}
 \boxed{
 \langle0|\Tord\{\psi_1\bar\psi_2\psi_3\bar\psi_4\}|0\rangle
 =S_{12}S_{34}-S_{14}S_{32},}
 \label{eq:fermion-wick-four-r68}
\end{equation}

\begin{proposition}[费米 Gaussian 态的四点配对]
设奇场满足 CAR，态为零均值准自由态；乘积顺序固定，配对符号由奇置换决定。则有以下结论（原式 \eqnrefs{eq:fermion-wick-four-r68}）：
\begin{equation*}
 \boxed{
 \langle0|\Tord\{\psi_1\bar\psi_2\psi_3\bar\psi_4\}|0\rangle
 =S_{12}S_{34}-S_{14}S_{32},}
\end{equation*}
\end{proposition}
\begin{proof}
正文\eqnrefs{eq:generic-fermion-contraction-dp8} 定义了自由费米场的一次收缩。对两对费米场只有直接配对和交叉配对两种完全收缩，而后者相对前者多一次 Grassmann 奇置换；把两项相加便得到正文\eqnrefs{eq:fermion-wick-four-r68}。



考虑固定初始顺序
\begin{equation*}
 \mathscr W_4
 =\langle0|T\{\psi_1\bar\psi_2\psi_3\bar\psi_4\}|0\rangle.
\end{equation*}
自由真空期望值必须把每个 $\psi$ 与一个 $\bar\psi$ 配对。第一种配对是 $(1,2)(3,4)$；两对在原顺序中已经相邻，因此不需要额外奇置换，贡献
\begin{equation*}
 +S_{12}S_{34}.
\end{equation*}
第二种配对是 $(1,4)(3,2)$。$\bar\psi_4$ 与 $\psi_1$ 收缩时必须越过中间的费米算符；等价地说，从配对 $(1,2)(3,4)$ 变到交叉配对改变了一次配对排列的宇称，因此贡献带负号。于是
\begin{equation*}
 {
 \mathscr W_4=S_{12}S_{34}-S_{14}S_{32}.}
 \end{equation*}
这个负号不是某个 QED 顶角的性质；即使没有电磁相互作用，自由费米高斯也已经具有它。

一般地，考虑 $n$ 个 $\psi_i$ 与 $n$ 个 $\bar\psi_j$。每个完全收缩等价于一个置换 $\pi\in S_n$，把第 $i$ 个 $\psi_i$ 配到 $\bar\psi_{\pi(i)}$。每次交换 Grassmann-odd 算符都产生 $-1$，故
\begin{equation*}
 \langle0|T\{\psi_1\bar\psi_1\cdots\psi_n\bar\psi_n\}|0\rangle
 =\sum_{\pi\in S_n}
 \operatorname{sgn}(\pi)
 \prod_{i=1}^{n}S_{i,\pi(i)}.
 \end{equation*}
右边正是矩阵 $S_{ij}$ 的行列式：
\begin{equation*}
 {
 \langle T\,\psi_1\bar\psi_1\cdots\psi_n\bar\psi_n\rangle_0
 =\det[S_{ij}].}
 \end{equation*}
与此相对，玻色高斯的完全收缩是不含 $\operatorname{sgn}(\pi)$ 的积和式配对和。Grassmann 高斯积分产生行列式，普通高斯积分产生行列式的负幂；它们与这里的 Wick 组合计数是同一反对易结构的不同表现。
\end{proof}

其中 $S_{ij}\equiv\langle0|\Tord\{\psi_i\bar\psi_j\}|0\rangle$。

在微扰展开中，若一组费米缩并首尾相接形成内部闭环，把该闭环上的 Grassmann 场重排为一个有序矩阵迹时会比对应的玻色配对多一次奇置换，因此每个独立闭合费米子圈带有
\begin{equation}
 \boxed{\text{each closed fermion loop contributes an overall factor }-1.}
 \label{eq:generic-closed-fermion-loop-r68}
\end{equation}

\begin{proposition}[闭费米收缩环的符号]
在 CAR 和固定时序的 Wick 展开中，比较具有同一外线顺序的完整闭合费米收缩链。则有以下结论（原式 \eqnrefs{eq:generic-closed-fermion-loop-r68}）：
\begin{equation*}
 \boxed{\text{each closed fermion loop contributes an overall factor }-1.}
\end{equation*}
\end{proposition}
\begin{proof}
在相互作用展开中，若若干费米收缩连成闭环，例如取如下示意性的双线性乘积

\begin{equation*}
 (\bar\psi_1\Gamma_1\psi_1)
 (\bar\psi_2\Gamma_2\psi_2)\cdots
 (\bar\psi_n\Gamma_n\psi_n),
\end{equation*}
并取循环收缩
\begin{equation*}
 \psi_1\leftrightarrow\bar\psi_2,
 \quad
 \psi_2\leftrightarrow\bar\psi_3,
 \quad\ldots\quad
 \psi_n\leftrightarrow\bar\psi_1,
\end{equation*}
则配对置换是一个长度为 $n$ 的循环。长度为 $n$ 的循环可分解成 $n-1$ 次对换，所以其宇称为 $(-1)^{n-1}$。另一方面，把每个双线性型中的 $\bar\psi_i$ 移到用于写矩阵迹的统一位置还产生 $(-1)^n$。两者相乘为
\begin{equation*}
 (-1)^{n-1}(-1)^n=-1.
\end{equation*}
因此一个闭合费米子圈相对于同样连接结构的对易的场多出
\begin{equation*}
 {-1.}
 \end{equation*}
\chapref{chap:electromagnetic-scattering}的 QED 真空极化的迹
$-\operatorname{Tr}[\gamma^\mu S_F\gamma^\nu S_F]$
就是这个一般结论的具体实例。
\end{proof}

具体相互作用顶角只决定环中的矩阵、内部量子数与传播子；这个整体负号在进入 QED、QCD 等具体理论之前已经由费米反对易代数固定。


为把 Wick 定理的组合计数连接到图规则，取无量纲耦合 $g_4$ 的标量相互作用
\begin{equation*}
 \Lag_{\rm int}=-\frac{g_4}{4!\,\hbar c}\varphi^4.
\end{equation*}
在 Dyson 展开的一阶，四点函数中与一个局域 $\varphi^4$ 插入对应的项是
\begin{equation}
 G^{(4,1)}(x_1,x_2,x_3,x_4)
 =-\frac{\ii g_4}{4!\,\hbar^2c^2}
 \int\dd^4 z\,
 \langle0|\Tord\{\varphi_1\varphi_2\varphi_3\varphi_4\varphi^4(z)\}|0\rangle .
 \label{eq:phi4-first-dyson-r68}
\end{equation}
只取把四个外场各自连接到 $z$ 点一个场的连通树级收缩，才得到 $g_4$ 的最低非平凡四点 Green 函数
\begin{equation}
 {
 G_c^{(4)}(x_1,x_2,x_3,x_4)
 =-\frac{\ii g_4}{\hbar^2c^2}
 \int\dd^4 z\,\prod_{a=1}^{4}\Delta_F(x_a-z).}
 \label{eq:phi4-tree-green-r3}
\end{equation}
$1/4!$ 与四条外部场插入同顶角上四个场的 $4!$ 种收缩正好抵消。Fourier 变换后，每个顶角的位置积分给出四动量守恒；未被这些 $\delta$ 函数固定的内部动量就是圈动量积分变量。

\begin{figure}[htbp]
 \centering
 \begin{tikzfig}[wfA]
  \coordinate (v) at (0,0);
  \draw (v)--(-2,1.5) node[left]{$x_1$};
  \draw (v)--(-2,-1.5) node[left]{$x_2$};
  \draw (v)--(2,1.5) node[right]{$x_3$};
  \draw (v)--(2,-1.5) node[right]{$x_4$};
  \fill (v) circle (2pt) node[above=5pt]{$z$};
 \end{tikzfig}
 \caption{$\varphi^4$ 理论四点 Green 函数的最低阶连通图。图只是\eqnrefs{eq:phi4-tree-green-r3} 中四个传播子与一次相互作用插入的图形缩写。}
 \label{fig:phi4-first-diagram-r3}
\end{figure}

同一 Dyson 展开在有限作用时间下还给出 Fermi 黄金律的来源。对时间无关微扰 $V$ 在时间窗 $T$ 内的一阶跃迁，定义能量失谐 $\Delta E\equiv E_f-E_i$，则有限时间概率核为
\begin{equation}
 {
 |c_f^{(1)}(T)|^2
 =\frac{|V_{fi}|^2T^2}{\hbar^2}
 \operatorname{sinc}^2\!\left(\frac{\Delta E\,T}{2\hbar}\right).}
 \label{eq:golden-sinc-r3}
\end{equation}
只有对平滑连续态密度 $\nu_f(E)$ 积分并取长时间极限，才可使用 $\sinc^2\to\delta$ 的分布极限，得到
\begin{equation}
 {
 w_{i\to f}
 =\frac{2\pi}{\hbar}|V_{fi}|^2\nu_f(E_i).}
 \label{eq:fermi-golden-r3}
\end{equation}
因此有限时间振幅先产生宽度约为 $1/T$ 的 sinc 谱核；在对平滑末态谱求和并取 $T\to\infty$ 后，该核才以分布极限收缩为能量守恒 $\delta$ 函数。

前面的自由场配对定理和图规则仍然使用自由场的产生、湮灭算符，但真实散射实验中的“入射电子”和“出射电子”属于相互作用理论。要把相关函数连接到实验散射振幅，必须先说明相互作用理论为什么在遥远过去和未来仍能定义自由的入射、出射粒子。

设相互作用理论存在稳定的单粒子态，物理质量为 $m_{\rm phys}$。把完整 Heisenberg 场记为 $\varphi_H(x)$。在遥远过去和遥远未来，彼此分离的稳定粒子已经离开有限的散射区，因此可分别引入自由的入场和出场
\begin{equation}
 \left(\Box+\frac{m_{\rm phys}^2c^2}{\hbar^2}\right)
 \varphi_{\rm in/out}(x)=0.
 \label{eq:asymptotic-free-field-dp48}
\end{equation}
它们各自具有与自由场相同形式的模展开，只是产生、湮灭算符分别记为 $a_{\rm in/out}(\vb{p})$。于是
\begin{equation*}
 |\vb{p};{\rm in}\rangle
 =a_{\rm in}^\dagger(\vb{p})|\Omega\rangle,
 \qquad
 |\vb{p};{\rm out}\rangle
 =a_{\rm out}^\dagger(\vb{p})|\Omega\rangle,
\end{equation*}
其中 $|\Omega\rangle$ 是相互作用理论的真空。多粒子入态、出态由相应产生算符继续构造，而实验散射振幅的起点正是
\begin{equation}
 {S_{fi}=\langle f,{\rm out}|i,{\rm in}\rangle.}
 \label{eq:in-out-smatrix-dp48}
\end{equation}
这一定义把“远处可以识别的粒子”与“相互作用区内复杂的量子场”明确分开。

对无导数标量相互作用，$\pi_H=c^{-2}\partial_t\varphi_H$。自由场模展开~\eqnrefs{eq:scalar-mode-r2} 可以逐模反演，因此对完整 Heisenberg 场定义有限时刻的单粒子投影算符
\begin{equation}
 \hat a_t(\vb{p})
 \equiv\int\dd^3 x\,\ee^{+\ii p\cdot x/\hbar}
 \left[
 \frac{1}{\hbar c}\sqrt{\frac{E_{\vb{p}}}{2}}\,\varphi_H(x)
 +\frac{\ii c}{\sqrt{2E_{\vb{p}}}}\,\pi_H(x)
 \right],
 \label{eq:lsz-finite-time-mode-projection-dp52}
\end{equation}
其中指数中的 $p^0=E_{\vb{p}}/c$ 取物理质量壳值。若 $\varphi_H$ 本身就是自由场，\eqnrefs{eq:lsz-finite-time-mode-projection-dp52} 精确返回 $\hat a(\vb{p})$；存在相互作用时，它会随 $t$ 改变，因为局域场在散射区内不断与多粒子扇区混合。

稳定单粒子极点意味着在足够早和足够晚的时刻，完整场在单粒子方向上的投影趋于自由入、出模。用极点留数 $Z$ 修正局域场与归一化粒子态之间的重叠，定义
\begin{equation}
 {
 \hat a_{\rm in/out}(\vb{p})
 =Z^{-1/2}\lim_{t\to\mp\infty}\hat a_t(\vb{p}).}
 \label{eq:asymptotic-annihilator-projection-dp52}
\end{equation}
对\eqnrefs{eq:lsz-finite-time-mode-projection-dp52} 求时间导数，再把空间 Laplace 项分部积分到平面波上，得到以下单腿恒等式。
\begin{proposition}[稳定标量单腿的约化恒等式]
设标量场具有物理质量 $m_{\rm phys}$、正极点留数 $Z>0$，并假设经波包涂抹的入、出单粒子模投影极限存在。场在所论矩阵元上可作相应时间微分，空间分部积分无边界通量，时间积分取与渐近极限一致的分布处方。则正文的有限时刻模投影满足
\begin{equation}
 {
 \hat a_{\rm out}(\vb{p})-\hat a_{\rm in}(\vb{p})
 =\frac{\ii Z^{-1/2}}{\sqrt{2E_{\vb{p}}}}
 \int\dd^4 x\,
 \ee^{+\ii p\cdot x/\hbar}
 \left(\Box+\frac{m_{\rm phys}^2c^2}{\hbar^2}\right)
 \varphi_H(x).}
 \label{eq:lsz-one-leg-position-dp52}
\end{equation}
\end{proposition}
\begin{proof}
先验证有限时刻投影式的归一化。把自由场展开~\eqnrefs{eq:scalar-mode-r2} 与

\begin{equation*}
 \pi(x)=\frac{1}{c^2}\partial_t\varphi(x)
 =-\ii\int\frac{\dd^3 q}{(2\pi\hbar)^3}
 \sqrt{\frac{E_{\vb{q}}}{2c^2}}
 \left[
 a(\vb{q})\ee^{-\ii q\cdot x/\hbar}
 -a^\dagger(\vb{q})\ee^{+\ii q\cdot x/\hbar}
 \right]
\end{equation*}
代入\eqnrefs{eq:lsz-finite-time-mode-projection-dp52}。对 $a(\vb{q})$ 分支，方括号中的两个系数分别给
\begin{align*}
 \frac{1}{\hbar c}\sqrt{\frac{E_{\vb{p}}}{2}}
 \frac{\hbar c}{\sqrt{2E_{\vb{q}}}}
 &\xrightarrow{\vb{q}=\vb{p}}\frac12,\\
 \frac{\ii c}{\sqrt{2E_{\vb{p}}}}
 \left(-\ii\sqrt{\frac{E_{\vb{q}}}{2c^2}}\right)
 &\xrightarrow{\vb{q}=\vb{p}}\frac12.
\end{align*}
空间积分给
\begin{equation*}
 \int\dd^3 x\,
 \ee^{-\ii(\vb{p}-\vb{q})\cdot\vb{x}/\hbar}
 =(2\pi\hbar)^3\delta^{(3)}(\vb{p}-\vb{q}),
\end{equation*}
因此两个 $1/2$ 相加后恰好得到 $a(\vb{p})$。对 $a^\dagger(\vb{q})$ 分支，空间积分要求 $\vb{q}=-\vb{p}$，而两个系数变成 $+1/2$ 与 $-1/2$，所以完全抵消。由此，\eqnrefs{eq:lsz-finite-time-mode-projection-dp52} 对自由场确实是模展开的精确逆变换。

现在把场换成完整 Heisenberg 场，并记
\begin{equation*}
 f_p(x)\equiv \ee^{+\ii p\cdot x/\hbar},
 \qquad p^0=E_{\vb{p}}/c.
\end{equation*}
用 $\pi_H=c^{-2}\partial_t\varphi_H$，\eqnrefs{eq:lsz-finite-time-mode-projection-dp52} 可写成
\begin{equation*}
 a_t(\vb{p})
 =\int\dd^3 x\,f_p(x)
 \left[
 A_p\varphi_H(x)+B_p\partial_t\varphi_H(x)
 \right],
\end{equation*}
其中
\begin{equation*}
 A_p=\frac{1}{\hbar c}\sqrt{\frac{E_{\vb{p}}}{2}},
 \qquad
 B_p=\frac{\ii}{c\sqrt{2E_{\vb{p}}}}.
\end{equation*}
对时间求导时，$\partial_t f_p=(\ii E_{\vb{p}}/\hbar)f_p$，所以
\begin{align*}
 \frac{\dd a_t}{\dd t}
 =\int\dd^3 x\,f_p
 \Bigl[&\frac{\ii E_{\vb{p}}}{\hbar}A_p\varphi_H
 +\left(A_p+\frac{\ii E_{\vb{p}}}{\hbar}B_p\right)
 \partial_t\varphi_H
 +B_p\partial_t^2\varphi_H\Bigr].
\end{align*}
中间的一阶时间导数系数为
\begin{align*}
 A_p+\frac{\ii E_{\vb{p}}}{\hbar}B_p
 &=\frac{1}{\hbar c}\sqrt{\frac{E_{\vb{p}}}{2}}
 -\frac{E_{\vb{p}}}{\hbar c\sqrt{2E_{\vb{p}}}}=0.
\end{align*}
余下两项可提出 $B_p$：
\begin{equation*}
 \frac{\dd a_t}{\dd t}
 =\frac{\ii}{c\sqrt{2E_{\vb{p}}}}
 \int\dd^3 x\,f_p
 \left[
 \partial_t^2+\frac{E_{\vb{p}}^2}{\hbar^2}
 \right]\varphi_H.
\end{equation*}
质量壳关系
\begin{equation*}
 E_{\vb{p}}^2=m_{\rm phys}^2c^4+c^2\vb{p}^2
\end{equation*}
把第二项分成质量项与空间动量项。空间部分做两次分部积分；若矩阵元中的波包在空间无穷远衰减，边界项消失，并有
\begin{equation*}
 \int\dd^3 x\,f_p\,(-c^2\nabla^2\varphi_H)
 =\int\dd^3 x\,(-c^2\nabla^2 f_p)\varphi_H
 =\frac{c^2\vb{p}^2}{\hbar^2}
 \int\dd^3 x\,f_p\varphi_H.
\end{equation*}
因此
\begin{equation*}
 \frac{\dd a_t}{\dd t}
 =\frac{\ii c}{\sqrt{2E_{\vb{p}}}}
 \int\dd^3 x\,f_p(x)
 \left(\Box+\frac{m_{\rm phys}^2c^2}{\hbar^2}\right)
 \varphi_H(x).
\end{equation*}
对 $t$ 从 $-\infty$ 积到 $+\infty$，并使用 $\dd^4 x=c\,\dd t\,\dd^3 x$，得到
\begin{equation*}
 \lim_{t\to+\infty}a_t(\vb{p})
 -\lim_{t\to-\infty}a_t(\vb{p})
 =\frac{\ii}{\sqrt{2E_{\vb{p}}}}
 \int\dd^4 x\,f_p(x)
 \left(\Box+\frac{m_{\rm phys}^2c^2}{\hbar^2}\right)
 \varphi_H(x).
\end{equation*}
最后用渐近定义~\eqnrefs{eq:asymptotic-annihilator-projection-dp52} 给两端同时乘 $Z^{-1/2}$，即得到\eqnrefs{eq:lsz-one-leg-position-dp52}。

该计算适用于自旋 $0$ 外腿：稳定极点固定单粒子留数，而 $t\to\pm\infty$ 的波包投影存在性已列为命题假设，空间分部积分要求波包或矩阵元在无穷远没有边界通量，而标量质量壳把 $E_{\vb{p}}^2$ 化成自由标量二次核 $\Box+m_{\rm phys}^2c^2/\hbar^2$。电子外腿对应一阶 Dirac 逆算符，必须使用费米场的 LSZ 约化。
\end{proof}
这里 $x^0=ct$，所以 $\dd^4 x=c\,\dd t\,\dd^3 x$。\eqnrefs{eq:lsz-one-leg-position-dp52} 已经建立了外腿约化的核心桥梁：每把一个标量渐近粒子从外态移入时序关联函数，就会在相应外部坐标上作用一次自由标量逆核 $\Box+m_{\rm phys}^2c^2/\hbar^2$。电子外腿不能沿用这一二阶算符；对自旋 $1/2$ 电子，截肢由一阶 Dirac 逆算符 $\ii\hbar\slashed\partial-m_{\mathrm e} c$ 完成。第 4 章将从 Dirac 二点函数本身建立其极点分子与外腿结构。对所有外腿分别作用与其 Lorentz 表示相匹配的自由逆算符，再 Fourier 变换，才得到相应的动量空间截肢公式。这整套由单粒子极点把时序相关函数转换成渐近散射振幅的步骤称为\emph{Lehmann--Symanzik--Zimmermann 约化}，简称 LSZ 约化。

完整场并不等于入场或出场。稳定单粒子态在精确二点函数中表现为与连续谱分支切割分离的孤立极点：
\begin{equation}
 {
 \widetilde G_2(p)
 =\frac{\ii\hbar^3c\,Z}{p^2-m_{\rm phys}^2c^2+\ii0^+}
 +\widetilde G_{2,\rm cont}(p).}
 \label{eq:lsz-pole-plus-cont-r7}
\end{equation}
$Z$ 是所选局域场与物理单粒子态的无量纲交叠留数；等价地说，$Z^{1/2}$ 衡量 $\varphi_H$ 在稳定单粒子极点附近与相应渐近自由场的重叠。因此 LSZ 中出现的 $Z^{-1/2}$ 不是人为补上的归一化常数，而是把“局域场插入”换成“归一化渐近粒子”的必要因子。

LSZ 外腿要求与连续谱分开的稳定单粒子极点。具有有限寿命的共振对应复能量平面上的极点，不能作为 $t\to\pm\infty$ 的自由渐近粒子；它只能通过内部共振传播或连续谱进入散射振幅。

对连通 $n$ 点函数，最清楚的写法不是用一个模糊的“近似于”符号，而是在每个外部质量壳附近作多变量 Laurent 分解。定义
\[
 \zeta_i\equiv p_i^2-m_i^2c^2,
\]
并先把 $\zeta_i$ 视为复变量；Feynman 的 $+\ii0^+$ 只规定最后从哪一侧取实轴边界值。则在稳定单粒子极点的邻域内可以写成精确分解
\begin{equation}
 {
 \widetilde G_{n,c}
 =
 (2\pi\hbar)^4\delta^{(4)}\!\left(\sum_i\sigma_i^{\rm io}p_i\right)
 \left[\prod_i
 \frac{\ii\hbar^3c\,Z_i^{1/2}}
 {\zeta_i+\ii0^+}\right]
 \ii\mathfrak A_{fi}(\{p_i\})
 +\widetilde R_n(\{p_i\}),}
 \label{eq:LSZ-factorization-r7}
\end{equation}
其中入射/出射外线分别取 $\sigma_i^{\rm io}=+1,-1$；$\mathfrak A_{fi}(\{p_i\})$ 在所有 $\zeta_i=0$ 的联合邻域内正则，而余项 $\widetilde R_n$ 的定义性质是它至少缺少一条外部单粒子极点，因此
\[
 \lim_{\zeta_1,\ldots,\zeta_n\to0}
 \left(\prod_i\zeta_i\right)\widetilde R_n=0.
\]
这一定义把“外腿极点”和“真正的相互作用核”严格分开。于是 LSZ 约化的核心就是先用逆极点因子截去外线传播子，再取在壳极限：
\begin{equation}
 {
 \lim_{p_i^2\to m_i^2c^2}
 \left[\prod_i
 Z_i^{-1/2}\frac{p_i^2-m_i^2c^2}{\ii\hbar^3c}\right]
 \widetilde G_{n,c}
 =(2\pi\hbar)^4\delta^{(4)}(P_f-P_i)\,\ii\mathfrak A_{fi}.}
 \label{eq:lsz-momentum-reduction-r7}
\end{equation}
对上面的 $\varphi^4$ 树图，四条外线传播子全部被这一操作消去，只留下
\begin{equation}
 {
 (2\pi\hbar)^4\delta^{(4)}(p_1+p_2-p_3-p_4)
 \left(-\frac{\ii g_4}{\hbar^2c^2}\right).}
 \label{eq:phi4-lsz-result-r7}
\end{equation}
因此外线传播子的截除来自精确单粒子极点的因子化。旋量与光子外线分别由相应的矩阵值逆动能算符完成截肢。

Dyson 级数产生自由场的时序乘积，Wick 定理把它们化为传播子与顶角，有限时间谱核说明能量守恒如何从长时间极限进入跃迁率，而精确单粒子极点再经 LSZ 约化把 Green 函数连接到渐近 $S$ 矩阵元。这四层共同构成微扰散射的算符基础。

% -----------------------------------------------------------------------------
% -----------------------------------------------------------------------------

\begin{exbox}[四点 Green 函数到截面]
取稳定的实标量粒子，使用树级留数 $Z=1$，相互作用为
\begin{equation*}
 \Lag_{\rm int}=-\frac{g_4}{4!\,\hbar c}\varphi^4,
 \qquad [g_4]=1,
\end{equation*}
物理外态的四动量满足 $p_i^2=m^2c^2$。从四点连通函数出发，用 LSZ 约化求在壳振幅，再计算二体散射的角分布与总截面。

\textbf{解。} 在动量空间中，四点连通函数的最低阶奇异部分为
\begin{equation}
 \begin{aligned}
 \widetilde G_{4,c}^{(1)}
 =&\,(2\pi\hbar)^4\delta^{(4)}(p_1+p_2-p_3-p_4)
 \left(-\frac{\ii g_4}{\hbar^2c^2}\right)\\
 &\times\prod_{a=1}^{4}
 \frac{\ii\hbar^3c}{p_a^2-m^2c^2+\ii0^+}.
 \end{aligned}
 \label{eq:phi4-fourpoint-poles-dp64}
\end{equation}
四维 $\delta$ 函数来自顶角位置积分，中间常数来自局域四价顶角，四个简单极点来自相互作用点到外部场插入点的传播。LSZ 约化把这些场插入转换成渐近单粒子外态，因此需要消去外腿极点。

对第 $a$ 条外腿定义逆极点因子
\begin{equation*}
 R_a(p_a)=\frac{p_a^2-m^2c^2}{\ii\hbar^3c}.
\end{equation*}
它与相应传播子在在壳留数极限下相消：
\[
 \lim_{p_a^2\to m^2c^2}
 R_a(p_a)\frac{\ii\hbar^3c}{p_a^2-m^2c^2+\ii0^+}=1.
\]
对四条外腿分别实施这一约化后，保留的顶角与总动量守恒因子为
\begin{equation*}
 (2\pi\hbar)^4\delta^{(4)}(P_f-P_i)
 \left(-\frac{\ii g_4}{\hbar^2c^2}\right).
\end{equation*}
按前文的国际单位制归一化约定，提出公共场归一化因子 $1/(\hbar^2c^2)$，余下系数满足 $\ii\mathcal M=-\ii g_4$，所以在壳约化振幅为
\begin{equation}
 \mathcal M_{\varphi\varphi\to\varphi\varphi}=-g_4.
 \label{eq:phi4-reduced-amplitude-dp64}
\end{equation}
外腿传播子已从振幅中消去，四个外粒子则由在壳动量标签标识。

将\eqnrefs{eq:phi4-reduced-amplitude-dp64} 代入\eqnrefs{eq:2to2-cross-section-master-si-r9}。等质量弹性散射满足 $|\vb{p}_f|=|\vb{p}_i|$，末态两个标量完全相同，因此全立体角计数需要除以 $2!$，得到
\begin{equation*}
 \frac{\dd\sigma}{\dd\Omega}
 =\frac12\frac{\hbar^2}{64\pi^2s}g_4^2
 =\frac{\hbar^2g_4^2}{128\pi^2s},
\end{equation*}
接触顶角没有 $t$ 或 $u$ 依赖，所以角分布各向同性。用 $\int\dd\Omega=4\pi$ 积分整个立体角，得到
\begin{equation}
 \sigma_{\rm tot}
 =\frac{\hbar^2g_4^2}{32\pi s}.
 \label{eq:phi4-total-cross-dp64}
\end{equation}
若用质心能量 $E_{\rm CM}=c\sqrt{s}$ 表示，便有
\begin{equation*}
 \sigma_{\rm tot}
 =\frac{g_4^2}{32\pi}
 \left(\frac{\hbar c}{E_{\rm CM}}\right)^2.
\end{equation*}
这里 $\hbar c/E_{\rm CM}$ 具有长度量纲，其平方给出截面的量纲。在这一树级接触模型中，固定 $g_4$ 时截面按 $E_{\rm CM}^{-2}$ 衰减。作为数量级估算，取 $g_4=0.50$、$E_{\rm CM}=1.00\,\mathrm{GeV}$，用 $\hbar c=0.19732698\,\mathrm{GeV\,fm}$ 得
\begin{align*}
 \sigma_{\rm tot}
 &=\frac{0.50^2}{32\pi}
 \left(\frac{0.19732698\,\mathrm{GeV\,fm}}{1.00\,\mathrm{GeV}}\right)^2,\\
 &\simeq9.68\times10^{-5}\,\mathrm{fm^2}
 =9.68\times10^{-35}\,\mathrm{m^2}
 \simeq0.968\,\mu\mathrm b.
\end{align*}
此数值来自同一振幅与相空间，不增加模型假设。全同末态因子、国际单位制归一化和树级截断都保留在计算中；加入圈修正后，振幅及其能量依赖还会改变。
\end{exbox}







% Detailed derivations gathered at the end of the owning structure.



\begin{derivation}[四点树级函数的收缩与组合因子]

正文\eqnrefs{eq:phi4-first-dyson-r68} 是 $\varphi^4$ 理论四点函数的一阶 Dyson 插入。连通树级项要求每个外场都与顶点 $z$ 上四个场中的一个且仅一个收缩，因此共有 $4!$ 种相同配对；它们与作用量中的 $1/4!$ 抵消后得到正文\eqnrefs{eq:phi4-tree-green-r3}。



四点函数的一阶 Dyson 项为
\begin{equation*}
 G^{(4,1)}(x_1,x_2,x_3,x_4)
 =-\frac{\ii g_4}{4!\,\hbar^2c^2}
 \int\dd^4 z\,
 \langle0|\Tord\{\varphi_1\varphi_2\varphi_3\varphi_4\varphi^4(z)\}|0\rangle .
\end{equation*}
要得到“一个顶点连接四个外点”的连通树图，每个外部场必须与 $z$ 点四个场中的一个且仅一个收缩。先给 $\varphi_1$ 选 $z$ 点场有 $4$ 种；再给 $\varphi_2$ 选剩余场有 $3$ 种；之后分别有 $2$ 种和 $1$ 种，所以
\begin{equation*}
 N_{\rm contr}=4\times3\times2\times1=4!.
\end{equation*}
每一种收缩都给出同一个核乘积
\begin{equation*}
 \Delta_F(x_1-z)\Delta_F(x_2-z)\Delta_F(x_3-z)\Delta_F(x_4-z).
\end{equation*}
因此 $4!$ 个相同项与相互作用中的 $1/4!$ 抵消，得到
\begin{equation*}
 G_c^{(4)}
 =-\frac{\ii g_4}{\hbar^2c^2}
 \int\dd^4 z\,\prod_{a=1}^4\Delta_F(x_a-z),
\end{equation*}
即正文\eqnrefs{eq:phi4-tree-green-r3}。

再把四条外线分别 Fourier 展开。四个传播子的指数中所有 $z$ 依赖合成
\begin{equation*}
 \exp\!\left[-\frac{\ii}{\hbar}
 (p_1+p_2-p_3-p_4)\cdot z\right].
\end{equation*}
于是顶角位置积分给
\begin{equation*}
 \int\dd^4 z\,
 \ee^{-\ii(p_1+p_2-p_3-p_4)\cdot z/\hbar}
 =(2\pi\hbar)^4\delta^{(4)}(p_1+p_2-p_3-p_4).
\end{equation*}
所以顶角四动量守恒正是局域相互作用对顶角位置积分后的 Fourier $\delta$ 函数。
\end{derivation}

\begin{derivation}[有限时间跃迁概率与黄金法则]
记 $\Delta E=E_f-E_i$。一阶振幅为

\begin{align*}
 c_f^{(1)}(T)
 &=-\frac{\ii}{\hbar}V_{fi}
 \int_{-T/2}^{T/2}\dd t\,\ee^{\ii\Delta E t/\hbar}\\
 &=-\frac{\ii}{\hbar}V_{fi}
 \left[\frac{\hbar}{\ii\Delta E}
 \ee^{\ii\Delta E t/\hbar}\right]_{-T/2}^{T/2}\\
 &=-\frac{2\ii V_{fi}}{\Delta E}
 \sin\!\left(\frac{\Delta E T}{2\hbar}\right).
\end{align*}
把 $\sinc x\equiv\sin x/x$ 提出，得到
\begin{equation*}
 |c_f^{(1)}(T)|^2
 =\frac{|V_{fi}|^2T^2}{\hbar^2}
 \sinc^2\!\left(\frac{\Delta E T}{2\hbar}\right),
\end{equation*}
即\eqnrefs{eq:golden-sinc-r3}。它在有限 $T$ 时不是 $\delta$ 函数；主峰宽度量级是 $\hbar/T$。

对连续末态求和，概率为
\begin{equation*}
 P_{i\to f}(T)
 =\int\dd E\,\nu_f(E)
 \frac{|V_{fi}(E)|^2T^2}{\hbar^2}
 \sinc^2\!\left[\frac{(E-E_i)T}{2\hbar}\right].
\end{equation*}
这里 $\nu_f(E)$ 是末态每单位能量的态密度，SI 量纲为 $\mathrm{J^{-1}}$。有限时间谱核与 Dirac $\delta$ 分布的对应由下列核固定：
\begin{equation*}
 K_T(\Delta E)
 \equiv
 \frac{T}{2\pi\hbar}
 \sinc^2\!\left(\frac{\Delta E\,T}{2\hbar}\right).
\end{equation*}
它具有能量倒数量纲。作变量代换 $x=\Delta E T/(2\hbar)$ 后，
\begin{align*}
 \int_{-\infty}^{+\infty}\dd(\Delta E)\,K_T(\Delta E)
 &=\frac{T}{2\pi\hbar}
 \frac{2\hbar}{T}
 \int_{-\infty}^{+\infty}\dd x\,\sinc^2x\\
 &=\frac1\pi\int_{-\infty}^{+\infty}\dd x\,\sinc^2x=1.
\end{align*}
因此 $K_T$ 是面积固定为一、宽度随 $T^{-1}$ 收缩的归一化谱核；对任意在 $E_i$ 附近平滑的测试函数 $F(E)$，$T\to\infty$ 时满足
\begin{equation*}
 \int\dd E\,F(E)K_T(E-E_i)\longrightarrow F(E_i),
\end{equation*}
也就是 $K_T(\Delta E)\to\delta(\Delta E)$ 的分布极限。
令
\begin{equation*}
F(E)\equiv |V_{fi}(E)|^2\nu_f(E).
\end{equation*}
由 $K_T$ 的定义，有限时间概率可以严格改写为
\begin{equation*}
P_{i\to f}(T)
=\frac{2\pi T}{\hbar}
\int\dd E\,F(E)K_T(E-E_i).
\end{equation*}
再定义有限时间余项
\begin{equation*}
\delta F_T
\equiv
\int\dd E\,[F(E)-F(E_i)]K_T(E-E_i),
\end{equation*}
利用 $\int K_T\dd E=1$ 得到精确等式
\begin{equation*}
P_{i\to f}(T)
=\frac{2\pi T}{\hbar}[F(E_i)+\delta F_T].
\end{equation*}
前面已经证明 $K_T\to\delta$ 的分布极限，所以对满足该近似恒等核收敛条件的连续 $F$，有 $\lim_{T\to\infty}\delta F_T=0$。因此长时间跃迁率严格由极限定义
\begin{equation*}
\lim_{T\to\infty}\frac{P_{i\to f}(T)}{T}
=\frac{2\pi}{\hbar}F(E_i).
\end{equation*}
也就是
\begin{equation*}
 w_{i\to f}=\frac{2\pi}{\hbar}|V_{fi}|^2\nu_f(E_i),
\end{equation*}
即 Fermi 黄金律~\eqnrefs{eq:fermi-golden-r3}。能量 $\delta$ 分布由有限时间谱核与连续谱卷积的长时间极限产生；有限 $T$ 时应保留具有有限宽度的 $\sinc^2$ 核。
\end{derivation}

\begin{derivation}[外腿极点怎样给出散射振幅]
先只处理一条外腿，因为多腿公式只是把同一个极点提取操作重复 $n$ 次。设 $\zeta=p^2-m_{\rm phys}^2c^2$，并把 $\zeta$ 暂时视为复变量。对时序关联函数中与坐标 $x$ 对应的那一个场插入完整态，稳定单粒子态给出
\begin{equation*}
 \langle\Omega|\varphi_H(0)|p,\lambda\rangle
 =Z^{1/2}\,\mathcal N_p,
\end{equation*}
其中 $\mathcal N_p$ 是由全书协变态归一化固定的运动学因子。把该单粒子贡献与其余所有中间态分开，Fourier 变换后必然具有形式
\begin{equation*}
 \widetilde G_{n,c}(p,\ldots)
 =\frac{\ii\hbar^3c\,Z^{1/2}}
 {\zeta}\,
 \widetilde F_{n-1}(p;\ldots)
 +\widetilde R^{(1)}_n(p,\ldots),
\end{equation*}
这里 $\widetilde F_{n-1}$ 是已经把这一条外腿换成归一化单粒子态后的矩阵元；$\widetilde R^{(1)}_n$ 包含连续谱和在这一外部通道中正则的项。稳定粒子极点与多粒子阈值分离，所以在 $\zeta=0$ 的一个小邻域内
\begin{equation*}
 \lim_{\zeta\to0}\zeta\,\widetilde R^{(1)}_n=0.
\end{equation*}
因此极点留数不是一句“外腿传播子可以删掉”，而是如下复分析极限：
\begin{align*}
 \lim_{\zeta\to0}
 Z^{-1/2}\frac{\zeta}{\ii\hbar^3c}
 \widetilde G_{n,c}
 &=\lim_{\zeta\to0}
 \left[
 \widetilde F_{n-1}
 +Z^{-1/2}\frac{\zeta}{\ii\hbar^3c}
 \widetilde R^{(1)}_n
 \right]\\
 &=\widetilde F_{n-1}(p_{\rm on};\ldots).
\end{align*}
Feynman 的 $+\ii0^+$ 是实动量边界值的处方；截肢本身应理解为先在复 $\zeta$ 平面取极点留数，再回到质量壳，因此不会出现把 $\zeta/(\zeta+\ii0^+)$ 当作普通数值比值的歧义。

对第二、第三直到第 $n$ 条外腿重复同一操作。若定义 $\zeta_i=p_i^2-m_i^2c^2$，则联合极点邻域中的多变量 Laurent 展开可写成
\begin{align*}
 \widetilde G_{n,c}
 ={}&(2\pi\hbar)^4\delta^{(4)}\!\left(\sum_i\sigma_i^{\rm io}p_i\right)
 \left[\prod_i
 \frac{\ii\hbar^3c\,Z_i^{1/2}}{\zeta_i}\right]
 \ii\mathfrak A_{fi}(\{p_i\})
 +\widetilde R_n,\\
 &\lim_{\zeta_1,\ldots,\zeta_n\to0}
 \left(\prod_i\zeta_i\right)\widetilde R_n=0.
\end{align*}
第二行正是“余项至少缺少一条完整外腿极点”的数学含义。于是
\begin{align*}
 &\lim_{\zeta_1,\ldots,\zeta_n\to0}
 \left[\prod_iZ_i^{-1/2}
 \frac{\zeta_i}{\ii\hbar^3c}\right]
 \widetilde G_{n,c}\\
 &\qquad=(2\pi\hbar)^4\delta^{(4)}\!\left(\sum_i\sigma_i^{\rm io}p_i\right)
 \ii\mathfrak A_{fi}(\{p_{i,\rm on}\}),
\end{align*}
这就得到正文的动量空间 LSZ 公式。注意截肢和在壳极限的顺序来自极点留数的定义：先乘 $\zeta_i$ 提取留数，再令 $\zeta_i\to0$；若反过来，原 Green 函数正好落在极点上。

现在把这一抽象结构用于树级 $\varphi^4$。一次 Dyson 插入给
\begin{equation*}
 -\frac{\ii}{\hbar c}
 \int\dd^4x\,
 \frac{g_4}{4!\,\hbar c}\,
 \varphi_I(x)^4.
\end{equation*}
把四个外部场分别与顶角处四个场收缩有 $4!$ 种双射，恰好消掉相互作用中的 $1/4!$。每条收缩给一个自由 Feynman 传播子，所以 Fourier 变换后的连通四点函数为
\begin{equation*}
 \widetilde G^{(1)}_{4,c}
 =-\frac{\ii g_4}{\hbar^2c^2}
 (2\pi\hbar)^4\delta^{(4)}(p_1+p_2-p_3-p_4)
 \prod_{a=1}^{4}
 \frac{\ii\hbar^3c}{p_a^2-m^2c^2+\ii0^+}.
\end{equation*}
在树级 $Z=1$。令 $\zeta_a=p_a^2-m^2c^2$，逐腿乘逆极点因子后，乘积中的四条外腿逐一消去：
\begin{align*}
 &\left[\prod_{a=1}^{4}
 \frac{\zeta_a}{\ii\hbar^3c}\right]
 \widetilde G^{(1)}_{4,c}\\
 &\quad=-\frac{\ii g_4}{\hbar^2c^2}
 (2\pi\hbar)^4\delta^{(4)}(p_1+p_2-p_3-p_4)
 \prod_{a=1}^{4}
 \frac{\zeta_a}{\zeta_a+\ii0^+}.
\end{align*}
按上面说明的极点留数顺序取极限，四个因子都贡献单位留数，于是得到正文\eqnrefs{eq:phi4-lsz-result-r7}。因此树级顶角规则不是另记的一条口诀：$4!$ 来自 Wick 收缩计数，四条传播子来自外部收缩，而 LSZ 的四个逆极点恰好把这些外腿传播子截去，只留下局域顶角系数。
\end{derivation}

